\chapter{วิธีการดำเนินการวิจัย}
\label{chapter:experiment}

ในโครงงานนี้ ผู้จัดทำได้พัฒนาระบบแฝดดิจิทัลสำหรับการจำลองสภาพการจราจรภายในมหาวิทยาลัย ซึ่งจะรับข้อมูลภาพวิดีโอจากถนนจริง ทำการตรวจจับยานพาหนะด้วยแบบจำลองตรวจจับวัตถุ YOLOv8 และติดตามวัตถุข้ามเฟรมด้วยอัลกอริทึม ByteTrack ข้อมูลที่ได้จะถูกนำมาวิเคราะห์เพื่อจำแนกสภาพการจราจรออกเป็น 4 สถานะ และส่งผลลัพธ์ไปแสดงบนแบบจำลองสามมิติอย่างต่อเนื่อง

ขั้นตอนการพัฒนาระบบในโครงงานนี้แบ่งออกเป็น 6 ขั้นตอนหลัก ตั้งแต่หัวข้อที่ \ref{sec:data-collection} ถึงหัวข้อที่ \ref{sec:visualization} จากนั้นหัวข้อที่ \ref{sec:experiment-design} จึงอธิบายการออกแบบการทดลองและตัวชี้วัดที่ใช้ประเมินระบบ อย่างไรก็ตาม ก่อนที่จะลงลึกถึงรายละเอียดของแต่ละขั้นตอน หัวข้อที่ \ref{sec:system-overview} และ \ref{sec:architecture} จะนำเสนอภาพรวมและสถาปัตยกรรมของระบบ เพื่อแสดงให้เห็นถึงความเชื่อมโยงและการส่งข้อมูลระหว่างโมดูลต่าง ๆ ก่อน

\section{ภาพรวมของระบบ}
\label{sec:system-overview}

ระบบที่พัฒนาขึ้นในโครงงานนี้มีเป้าหมายเพื่อประมวลผลข้อมูลภาพวิดีโอจากกล้องที่ติดตั้งบนสะพานลอยคนข้ามเหนือถนน ให้กลายเป็นข้อมูลเชิงตัวเลขที่สามารถอธิบายสภาพการจราจรของถนนในช่วงนั้นได้ และนำข้อมูลดังกล่าวไปแสดงผลบนแบบจำลองสามมิติที่เปลี่ยนแปลงตามสภาพความเป็นจริงของถนนอย่างต่อเนื่อง

ข้อมูลที่ระบบรับเข้ามาคือไฟล์วิดีโอที่บันทึกจากถนนฉลองกรุง ส่วนข้อมูลที่ระบบส่งออกมีสองชุด ชุดแรกคือรายการยานพาหนะที่ข้ามเส้นที่ใช้นับ ซึ่งระบุชนิดและทิศทางของยานพาหนะแต่ละคันที่ข้ามเส้นไป ชุดที่สองคือคำตัดสินสภาพการจราจรของแต่ละช่วงเวลา ซึ่งจำแนกออกเป็นสี่สถานะ ได้แก่ การจราจรปกติ การจราจรหนาแน่น การจราจรติดขัดเคลื่อนตัวช้า และการจราจรติดขัดหยุดนิ่ง

เส้นทางของข้อมูลผ่านสามโมดูลที่ทำงานแยกกระบวนการกัน ดังแสดงในรูปที่ \ref{fig:system-architecture} โมดูลปัญญาประดิษฐ์ (AI Worker) จะอ่านภาพจากไฟล์วิดีโอ ตรวจจับและติดตามยานพาหนะ นับยานพาหนะที่ข้ามเส้นนับ และสรุปสภาพการจราจรเป็นช่วงเวลา จากนั้นส่งผลลัพธ์ไปยังโมดูลระบบหลังบ้าน (Backend) ซึ่งทำหน้าที่ตรวจสอบความถูกต้องของข้อมูลและกระจายต่อให้ผู้รับทุกราย ปลายทางคือโมดูลแสดงผล (Unity) ซึ่งสร้างยานพาหนะและปรับสถานะของฉากสามมิติตามข้อมูลที่ได้รับ รายละเอียดของแต่ละโมดูลจะอยู่ในหัวข้อที่ \ref{sec:architecture}

\begin{figure}[ht]
    \centering
    \includegraphics[width=140mm]{fig_3_1.png}
    \caption{ภาพรวมสถาปัตยกรรมของระบบ}
    \label{fig:system-architecture}
\end{figure}

เมื่อพิจารณาตามองค์ประกอบสามส่วนของแฝดดิจิทัลที่กล่าวไว้ในบทที่ \ref{chapter:literature-review} ฝั่งกายภาพของระบบนี้คือถนนจริงและกล้องที่บันทึกภาพ ฝั่งเสมือนคือแบบจำลองสามมิติ และการเชื่อมโยงข้อมูลคือโมดูลปัญญาประดิษฐ์ร่วมกับโมดูลระบบหลังบ้าน อย่างไรก็ตาม ข้อมูลไหลในทิศทางเดียวจากฝั่งกายภาพไปยังฝั่งเสมือนเท่านั้น ฝั่งเสมือนไม่สามารถส่งคำสั่งย้อนกลับไปควบคุมสิ่งใดในโลกจริงได้ ระบบนี้จึงอยู่ในระดับเงาดิจิทัล (Digital Shadow) ยังไม่ถึงระดับแฝดดิจิทัลเต็มรูปแบบ การยกระดับไปสู่การป้อนกลับสองทางถูกจัดไว้เป็นงานในอนาคต

\section{สถาปัตยกรรมของระบบ}
\label{sec:architecture}

\subsection{องค์ประกอบของระบบ}
\label{subsec:components}

ระบบประกอบด้วยสามโมดูลที่ทำงานแยกออกจากกัน และสื่อสารกันผ่านโปรโตคอลมาตรฐานคือ HTTP และเว็บซ็อกเก็ต การแยกเช่นนี้ทำให้แต่ละโมดูลพัฒนาและทดสอบได้อย่างอิสระ และทำให้สามารถแทนที่โมดูลใดโมดูลหนึ่งด้วยตัวจำลองในระหว่างการพัฒนาได้โดยที่โมดูลอื่นไม่รับรู้ความแตกต่าง

โมดูลปัญญาประดิษฐ์ (AI Worker) พัฒนาด้วยภาษาไพทอน ทำหน้าที่อ่านภาพจากไฟล์วิดีโอทีละเฟรม ตรวจจับและติดตามยานพาหนะ ตรวจการข้ามเส้นนับ คำนวณค่าความหนาแน่นและอัตราการไหลต่อช่วงเวลา ตัดสินสภาพการจราจร และส่งผลลัพธ์ออกไปยังโมดูลระบบหลังบ้าน โมดูลนี้ไม่จัดเก็บข้อมูลถาวร ไม่จัดการการเชื่อมต่อของผู้รับ และไม่รับรู้ว่าปลายทางของข้อมูลคือโมดูลใด

โมดูลระบบหลังบ้าน (Backend) พัฒนาด้วย Node.js ร่วมกับเฟรมเวิร์ก Express ทำหน้าที่รับข้อมูลผ่านจุดปลายทาง HTTP ตรวจสอบความถูกต้องของข้อมูลก่อนประมวลผลใด ๆ ปรับปรุงตัวนับสะสมในหน่วยความจำ และกระจายข้อมูลต่อไปยังผู้รับทุกรายผ่านเว็บซ็อกเก็ตด้วยไลบรารี Socket.IO โมดูลนี้ไม่ประมวลผลภาพ ไม่เรียกใช้แบบจำลองปัญญาประดิษฐ์ และไม่รับรู้ตรรกะการแสดงผลของโมดูลปลายทาง นอกจากนี้ยังมีหน้าเว็บสำหรับตรวจสอบระหว่างการพัฒนา ซึ่งรับข้อมูลชุดเดียวกันเพื่อให้ผู้พัฒนาเห็นข้อมูลที่ไหลผ่านระบบได้โดยไม่ต้องเปิดโมดูลแสดงผล

โมดูลแสดงผล (Unity) พัฒนาด้วยเกมเอนจิน Unity และภาษา C\# ทำหน้าที่เชื่อมต่อเว็บซ็อกเก็ตกับโมดูลระบบหลังบ้านเมื่อเริ่มทำงาน รับรายการยานพาหนะที่ข้ามเส้นนับ และสร้างยานพาหนะในฉากสามมิติตามชนิดและทิศทางที่ได้รับ การจัดการวัตถุใช้เทคนิคการนำวัตถุกลับมาใช้ซ้ำ (Object Pooling) โดยเมื่อยานพาหนะเคลื่อนที่จนสุดเส้นทาง วัตถุนั้นจะถูกปิดการทำงานและเก็บคืนเข้าคิวของแบบจำลองชนิดเดียวกัน เพื่อนำกลับมาใช้ใหม่ในครั้งถัดไป แทนการทำลายและสร้างวัตถุขึ้นใหม่ ซึ่งเป็นสาเหตุหลักของการกระตุกในเกมเอนจิน

ในเวอร์ชันต้นแบบนี้ โมดูลแสดงผลรับข้อมูลสองช่องทาง ได้แก่ รายการยานพาหนะที่ข้ามเส้นนับ (\texttt{spawn\_vehicle}) ซึ่งใช้สร้างยานพาหนะในฉาก และคำตัดสินสภาพการจราจรรายโซน (\texttt{traffic\_state}) ซึ่งใช้ปรับช่วงความเร็วของยานพาหนะในแต่ละโซนให้สอดคล้องกับสภาพการจราจร คำตัดสินเดียวกันนี้ยังตรวจสอบได้ผ่านหน้าเว็บสำหรับผู้พัฒนาด้วย นอกจากนี้ระบบยังไม่มีการจัดเก็บข้อมูลลงฐานข้อมูล ตัวนับสะสมทั้งหมดอยู่ในหน่วยความจำของโมดูลระบบหลังบ้านและจะถูกตั้งค่ากลับเป็นศูนย์เมื่อเริ่มต้นกระบวนการใหม่

\subsection{การไหลของข้อมูล}
\label{subsec:dataflow}

การไหลของข้อมูลตั้งแต่ไฟล์วิดีโอจนถึงการแสดงผล แสดงในรูปที่ \ref{fig:data-flow} และมีลำดับขั้นดังนี้

\begin{enumerate}
    \item โมดูลปัญญาประดิษฐ์อ่านภาพจากไฟล์วิดีโอทีละเฟรม ตรวจจับยานพาหนะด้วยแบบจำลอง YOLOv8 และติดตามยานพาหนะข้ามเฟรมด้วยตัวติดตาม ByteTrack เพื่อให้ยานพาหนะคันเดียวกันมีรหัสประจำตัวคงเดิม
    \item เมื่อยานพาหนะคันหนึ่งข้ามเส้นนับในทิศทางที่กำหนด ระบบสร้างเหตุการณ์การข้ามเส้นนับหนึ่งรายการ ซึ่งระบุรหัสประจำตัว ชนิด และทิศทางของยานพาหนะคันนั้น แล้วส่งด้วยคำสั่ง HTTP POST ไปยังโมดูลระบบหลังบ้าน
    \item ในขณะเดียวกัน ทุกๆช่วงเวลา 20 วินาที ระบบสรุปค่าความหนาแน่นและอัตราการไหลของแต่ละโซน ตัดสินสภาพการจราจรของช่วงเวลานั้น แล้วส่งไปยังโมดูลระบบหลังบ้านอีกจุดปลายทางหนึ่ง ข้อมูลสองสายนี้มีความถี่และความหมายต่างกัน สายแรกอธิบายยานพาหนะรายคัน ส่วนสายที่สองอธิบายสภาพของถนนโดยรวม จึงแยกช่องทางกันตั้งแต่ต้นทาง
    \item โมดูลระบบหลังบ้านตรวจสอบความถูกต้องของข้อมูลที่รับเข้ามาก่อนประมวลผลใด ๆ หากข้อมูลไม่ครบตามที่กำหนด หรือมีค่าที่อยู่นอกช่วงที่ยอมรับได้ โมดูลจะตอบกลับด้วยรหัสสถานะ HTTP 400 พร้อมระบุเหตุผลที่ปฏิเสธ และไม่กระจายข้อมูลนั้นต่อ การตรวจสอบที่ด่านนี้ทำให้ข้อมูลที่ผิดพลาดไม่ถูกส่งต่อไปสะสมเป็นความคลาดเคลื่อนในโมดูลปลายทาง และทำให้ระบุได้ทันทีว่าความผิดพลาดเกิดที่โมดูลต้นทาง
    \item หากข้อมูลผ่านการตรวจสอบ โมดูลระบบหลังบ้านปรับปรุงตัวนับสะสมแยกตามชนิดของยานพาหนะ ทิศทาง และรหัสกล้อง
    \item โมดูลระบบหลังบ้านกระจายข้อมูลต่อผ่านเว็บซ็อกเก็ตโดยไม่แก้ไขค่าใด ๆ ยกเว้นช่องทางที่ส่งให้โมดูลแสดงผล ซึ่งตัดฟิลด์ที่โมดูลนั้นไม่ได้ใช้ออกสี่ฟิลด์ ได้แก่ รุ่นของสัญญาข้อมูล เวลาในวิดีโอต้นทาง ลำดับเฟรม และค่าความเชื่อมั่นของการตรวจจับซึ่งถูกนำไปใช้ตั้งแต่ขั้นตรวจสอบและนับแล้ว การกำหนดให้โมดูลกลางไม่แปลงค่าข้อมูลเป็นหลักการสำคัญของระบบ เพราะทำให้ข้อมูลที่ปลายทางได้รับตรงกับข้อมูลที่ต้นทางส่งออกมาเสมอ และเมื่อผลลัพธ์ผิดพลาดจึงตรวจสอบย้อนกลับได้ว่าเกิดขึ้นที่โมดูลใด
    \item โมดูลแสดงผลรับข้อมูลจากช่องทางรายการยานพาหนะ แล้วนำโมเดลยานพาหนะออกมาแสดงผลในฉากตามชนิดและทิศทางที่ได้รับ ส่วนช่องทางสภาพการจราจรใช้สำหรับปรับการแสดงสถานะของฉากตามคำตัดสินที่ได้รับในแต่ละช่วงเวลา
\end{enumerate}

\begin{figure}[htbp]
    \centering
    \includegraphics[width=\textwidth, height=0.8\textheight, keepaspectratio]{fig_3_2.png}
    \caption{การไหลของข้อมูลตั้งแต่ไฟล์วิดีโอจนถึงการแสดงผล}
    \label{fig:data-flow}
\end{figure}

\subsection{สัญญาข้อมูลระหว่างโมดูล}
\label{subsec:data-contract}

เนื่องจากทั้งสามโมดูลพัฒนาด้วยภาษาต่างกันและทำงานแยกกระบวนการกัน ระบบจึงกำหนดโครงสร้างข้อมูลที่ใช้สื่อสารระหว่างกันไว้เป็นลายลักษณ์อักษร เรียกว่าสัญญาข้อมูล (Data Contract) โดยข้อมูลทุกชุดอยู่ในรูปแบบ JSON และมีฟิลด์ระบุรุ่นของสัญญากำกับไว้ เพื่อให้โมดูลปลายทางตรวจสอบได้ว่ากำลังรับข้อมูลรุ่นใด โมดูลระบบหลังบ้านทำหน้าที่บังคับใช้สัญญานี้ โดยตรวจสอบข้อมูลทุกชุดก่อนกระจายต่อ

สัญญาข้อมูลชุดแรกคือเหตุการณ์การข้ามเส้นนับ ส่งผ่านจุดปลายทาง \texttt{POST /api/ingest} หนึ่งรายการต่อยานพาหนะหนึ่งคันที่ข้ามเส้น มีรายละเอียดดังตารางที่ \ref{tab:spawn-contract}

\begin{table}[ht]
    \caption{สัญญาข้อมูลของเหตุการณ์การข้ามเส้นนับ}
    \label{tab:spawn-contract}
    \centering
    \footnotesize
    \begin{tabular}{|>{\raggedright\arraybackslash}p{30mm}|l|c|>{\raggedright\arraybackslash}p{62mm}|}
        \hline
        \textbf{ฟิลด์} & \textbf{ชนิดข้อมูล} & \textbf{บังคับ} & \multicolumn{1}{c|}{\textbf{ความหมายและเงื่อนไขการตรวจสอบ}} \\
        \hline
        \texttt{schema} & สายอักขระ & ไม่ & รุ่นของสัญญาข้อมูล ปัจจุบันคือ \texttt{spawn-event/0.2-draft} \\
        \hline
        \texttt{timestamp} & สายอักขระ & ใช่ & เวลาที่เกิดเหตุการณ์ ต้องอยู่ในรูปแบบ ISO 8601 \\
        \hline
        \texttt{cameraId} & สายอักขระ & ใช่ & รหัสกล้องต้นทาง ต้องไม่เป็นค่าว่าง \\
        \hline
        \texttt{videoTimeSec} & จำนวนจริง & ไม่ & เวลาในวิดีโอต้นทาง หน่วยเป็นวินาที ต้องไม่ติดลบ \\
        \hline
        \texttt{frameCount} & จำนวนเต็ม & ไม่ & ลำดับเฟรมในวิดีโอต้นทาง ต้องไม่ติดลบ \\
        \hline
        \texttt{trackId} & สายอักขระ & ใช่ & รหัสประจำตัวยานพาหนะ รูปแบบคือชนิดตามด้วยเลขสี่หลัก ต้องไม่เป็นค่าว่าง \\
        \hline
        \texttt{type} & สายอักขระ & ใช่ & ชนิดยานพาหนะ ต้องเป็นหนึ่งใน \texttt{car}, \texttt{truck}, \texttt{motorcycle}, \texttt{bus} \\
        \hline
        \texttt{direction} & สายอักขระ & ใช่ & ทิศทางการข้ามเส้น ต้องเป็น \texttt{in} หรือ \texttt{out} \\
        \hline
        \texttt{confidence} & จำนวนจริง & ใช่ & ค่าความเชื่อมั่นของการตรวจจับ ต้องอยู่ระหว่าง 0.0 ถึง 1.0 \\
        \hline
    \end{tabular}
\end{table}

สัญญาข้อมูลชุดที่สองคือคำตัดสินสภาพการจราจรต่อช่วงเวลา ส่งผ่านจุดปลายทาง \texttt{POST /api/traffic-state} หนึ่งรายการต่อหนึ่งช่วงเวลา มีรายละเอียดดังตารางที่ \ref{tab:trafficstate-contract}

\begin{table}[ht]
    \caption{ข้อกำหนดโครงสร้างข้อมูล (Data Contract) สำหรับผลการประเมินสภาพการจราจร}
    \label{tab:trafficstate-contract}
    \centering
    \footnotesize
    \begin{tabular}{|>{\raggedright\arraybackslash}p{42mm}|l|c|>{\raggedright\arraybackslash}p{58mm}|}
        \hline
        \textbf{ฟิลด์} & \textbf{ชนิดข้อมูล} & \textbf{บังคับ} & \multicolumn{1}{c|}{\textbf{ความหมายและเงื่อนไขการตรวจสอบ}} \\
        \hline
        \texttt{schema} & สายอักขระ & ไม่ & เวอร์ชันของโครงสร้างข้อมูล (ปัจจุบันคือ \texttt{traffic-state/0.1-draft}) \\
        \hline
        \texttt{timestamp} & สายอักขระ & ใช่ & เวลาที่ระบบประมวลผลเสร็จสิ้น (ต้องระบุในรูปแบบมาตรฐาน ISO 8601) \\
        \hline
        \texttt{cameraId} & สายอักขระ & ใช่ & รหัสประจำตัวกล้องที่ใช้ตรวจจับ (ห้ามเป็นค่าว่าง) \\
        \hline
        \texttt{windowStartSec} & จำนวนจริง & ใช่ & เวลาเริ่มต้นของการวัดข้อมูล (หน่วยเป็นวินาที) โดยค่าต้องไม่น้อยกว่าศูนย์ \\
        \hline
        \texttt{windowEndSec} & จำนวนจริง & ใช่ & เวลาสิ้นสุดของการวัดข้อมูล (หน่วยเป็นวินาที) โดยค่าต้องไม่น้อยกว่าศูนย์ และต้องมีค่ามากกว่า \texttt{windowStartSec} \\
        \hline
        \texttt{zones} & ออบเจกต์ & ใช่ & กลุ่มข้อมูลพื้นที่ตรวจวัด โดยใช้ชื่อโซน (\texttt{in} หรือ \texttt{out}) เป็นคีย์ ต้องมีอย่างน้อย 1 พื้นที่ ถนนทางเดียวมีเพียงโซนเดียวได้ \\
        \hline
        \texttt{zones.*.occupancy} & จำนวนจริง & ใช่ & ปริมาณยานพาหนะเฉลี่ยในพื้นที่ตรวจวัดต่อเฟรมภาพ โดยค่าต้องไม่น้อยกว่าศูนย์ \\
        \hline
        \texttt{zones.*.vehicleFlowRate} & จำนวนจริง & ใช่ & อัตราการไหลของยานพาหนะทั่วไป (รถยนต์ รถบรรทุก รถโดยสาร) มีหน่วยเป็นคัน/นาที โดยค่าต้องไม่น้อยกว่าศูนย์ \\
        \hline
        \texttt{zones.*.motorcycleFlowRate} & จำนวนจริง & ไม่ & อัตราการไหลของรถจักรยานยนต์ มีหน่วยเป็นคัน/นาที หากมีการระบุข้อมูล ค่าต้องไม่น้อยกว่าศูนย์ หากไม่ระบุให้ถือว่าเป็นศูนย์ \\
        \hline
        \texttt{zones.*.trafficState} & สายอักขระ & ไม่ & สถานะการจราจรเฉพาะของโซนนั้น ใช้ชุดค่าเดียวกับ \texttt{trafficState} ระดับบนสุด หากไม่ระบุ ให้โมดูลปลายทางใช้ค่าระดับบนสุดแทน \\
        \hline
        \texttt{trafficState} & สายอักขระ & ใช่ & สถานะการจราจรโดยรวม คือสถานะของโซนที่รุนแรงที่สุด โดยต้องระบุเป็นค่าใดค่าหนึ่งดังนี้: \texttt{normal}, \texttt{high\_density}, \texttt{slow\_moving}, หรือ \texttt{standstill} \\
        \hline
    \end{tabular}
\end{table}

ประเด็นสำคัญของสัญญาชุดที่สองอยู่ที่การแยกอัตราการไหลออกเป็นสองฟิลด์ กล่าวคือ \texttt{vehicleFlowRate} นับเฉพาะรถยนต์ รถบรรทุก และรถโดยสาร ส่วนรถจักรยานยนต์ถูกแยกไปไว้ที่ \texttt{motorcycleFlowRate} ต่างหาก การแยกตั้งแต่ระดับสัญญาข้อมูลเช่นนี้ทำให้ค่าที่ใช้ตัดสินสภาพการจราจรไม่ถูกรบกวนจากรถจักรยานยนต์ที่ยังเบียดตัวผ่านไปได้ในขณะที่ยานพาหนะประเภทอื่นหยุดนิ่ง ฟิลด์ \texttt{motorcycleFlowRate} กำหนดให้ไม่บังคับ เพื่อให้โมดูลระบบหลังบ้านยังรับข้อมูลจากโมดูลปัญญาประดิษฐ์รุ่นก่อนหน้าที่ยังไม่ได้แยกฟิลด์นี้ได้

เมื่อข้อมูลผ่านการตรวจสอบแล้ว โมดูลระบบหลังบ้านกระจายต่อผ่านเว็บซ็อกเก็ตสามช่องทาง ได้แก่ ช่องทาง \texttt{spawn} ซึ่งส่งข้อมูลเหตุการณ์การข้ามเส้นนับครบทุกฟิลด์ ช่องทาง \texttt{spawn\_vehicle} ซึ่งส่งข้อมูลชุดเดียวกันแต่ตัดเหลือเฉพาะห้าฟิลด์ที่โมดูลแสดงผลใช้จริง ได้แก่ \texttt{trackId}, \texttt{type}, \texttt{direction}, \texttt{cameraId} และ \texttt{timestamp} และช่องทาง \texttt{traffic\_state} ซึ่งส่งคำตัดสินสภาพการจราจรครบทุกฟิลด์

\section{การเก็บวิดีโอและการกำหนดพื้นที่ตรวจวัด}
\label{sec:data-collection}

\subsection{การสำรวจพื้นที่และการบันทึกวิดีโอ}
\label{subsec:site-survey}

พื้นที่ที่ใช้ในการศึกษาคือถนนฉลองกรุง ซึ่งเป็นถนนหลายช่องจราจรที่มีรถจักรยานยนต์เป็นสัดส่วนสูง ตรงกับลักษณะกระแสจราจรแบบผสมที่กล่าวไว้ในหัวข้อที่ \ref{subsec:mixed-traffic} และอยู่ในบริเวณที่เข้าถึงได้สำหรับการบันทึกภาพ

กล้องถูกติดตั้งบนสะพานลอยคนข้าม โดยหันกล้องไปตามแนวถนน มิใช่ตั้งฉากกับถนน  การเลือกมุมกล้องเช่นนี้มีเหตุผลสามประการ ประการแรกคือมุมมองตามแนวถนนทำให้เห็นถนนได้ยาวกว่า จึงวางโซนสังเกตการณ์ที่ครอบคลุมพื้นที่มากพอสำหรับการวัดความหนาแน่นได้ ประการที่สองคือความสูงของสะพานลอยช่วยลดการบดบังกันเองของยานพาหนะ ซึ่งเป็นปัญหาสำคัญเมื่อการจราจรหนาแน่น ประการที่สามคือมุมนี้ทำให้ยานพาหนะที่วิ่งเข้าหากล้องและวิ่งออกจากกล้องแยกออกจากกันได้ชัดเจนในภาพ

จากมุมกล้องดังกล่าว ผู้จัดทำกำหนดนิยามการเคลื่อนที่ในภาพไว้ว่า ยานพาหนะที่เคลื่อนที่
เข้าหากล้องจะปรากฏเป็นการเลื่อนลงด้านล่างของภาพ กำหนดให้เป็นค่า \texttt{toward}
ส่วนยานพาหนะที่เคลื่อนที่ออกจากกล้องจะปรากฏเป็นการเลื่อนขึ้นด้านบนของภาพ กำหนดให้เป็นค่า
\texttt{away} นิยามนี้เป็นเรื่องของเรขาคณิตในภาพล้วน ๆ และถูกใช้ในการตรวจทิศทาง
การข้ามเส้นตามหัวข้อที่ \ref{subsec:direction-check} เพื่อปฏิเสธยานพาหนะที่เคลื่อนที่
สวนทางกับที่โซนนั้นรับผิดชอบ

ส่วนทิศทางที่ระบบส่งออกไปยังโมดูลปลายทาง คือค่า \texttt{in} หรือ \texttt{out} นั้น
กำหนดจากชื่อของโซนโดยตรง ไม่ได้คำนวณจากทิศการเคลื่อนที่ในภาพ ผู้จัดทำแยกสองสิ่งนี้
ออกจากกันเนื่องจากความสัมพันธ์ระหว่างฝั่งถนนกับทิศการเคลื่อนที่ในภาพขึ้นอยู่กับว่ากล้องถูก
ติดตั้งหันไปทางใด การผูกค่าทั้งสองเข้าด้วยกันจึงทำให้ต้องแก้โปรแกรมทุกครั้งที่เปลี่ยนมุมกล้อง
ผลของการแยกคือโซนหนึ่งสามารถมีชื่อเป็น \texttt{in} ขณะที่ยานพาหนะในโซนนั้นเคลื่อนที่
ออกจากกล้อง (\texttt{away}) ได้ ซึ่งเป็นกรณีที่เกิดขึ้นจริงกับชุดข้อมูลที่ใช้ทดสอบ

อุปกรณ์ที่ใช้บันทึกคือโทรศัพท์เคลื่อนที่ Samsung รุ่น A54 ระบบปฏิบัติการ Android 16 บันทึกที่ความละเอียด 3840$\times$2160 จุดภาพ อัตราเฟรม 30 เฟรมต่อวินาที ติดตั้งบนสะพานลอยคนข้ามที่ความสูงประมาณ 5.5 เมตรจากผิวถนน ทำมุมก้มประมาณ 60 องศา บันทึกในช่วงเวลา 13 สิงหาคม พ.ศ.2569 ระหว่างเวลา 14.00 - 15.00 น.ภายใต้สภาพแสงกลางวันและอากาศแจ่มใส

การบันทึกทั้งหมดดำเนินการในสภาพแสงกลางวัน ระบบจึงยังไม่ได้รับการทดสอบในเวลากลางคืนหรือขณะฝนตก 

\begin{figure}[ht]
    \centering
    \includegraphics[width=130mm]{fig_3_3.png}
    \caption{ตำแหน่งกล้องและพื้นที่ที่บันทึกภาพ}
    \label{fig:camera-site}
\end{figure}

\subsection{การกำหนดโซนและเส้นนับ}
\label{subsec:zone-setup}

ก่อนที่ระบบจะประมวลผลวิดีโอได้ ผู้จัดทำต้องกำหนดพื้นที่ตรวจวัดบนภาพเสียก่อน โดยแต่ละโซนประกอบด้วยองค์ประกอบทางเรขาคณิตสามส่วนและทิศทางที่คาดหวังหนึ่งค่า ดังนี้

\begin{itemize}
    \item \textbf{โซนนับ} (\texttt{polygon}) เป็นรูปหลายเหลี่ยมที่ครอบช่องจราจรที่ต้องการนับ ทำหน้าที่เป็นตัวกรองเชิงพื้นที่ กล่าวคือ การข้ามเส้นนับจะถูกยอมรับก็ต่อเมื่อเกิดขึ้นภายในรูปหลายเหลี่ยมนี้เท่านั้น  การมีตัวกรองนี้ป้องกันไม่ให้ระบบนับยานพาหนะที่วิ่งอยู่บนถนนฝั่งตรงข้าม หรือวัตถุที่เคลื่อนที่อยู่นอกพื้นที่ที่สนใจ
    \item \textbf{เส้นนับ} (\texttt{line}) เป็นส่วนของเส้นตรงที่พาดขวางช่องจราจร ระบบจะนับยานพาหนะหนึ่งคันเมื่อจุดอ้างอิงของยานพาหนะนั้นเคลื่อนผ่านจากด้านหนึ่งของเส้นไปยังอีกด้านหนึ่ง รายละเอียดของการตรวจการข้ามเส้นอยู่ในหัวข้อที่ \ref{subsec:line-crossing}
    \item \textbf{โซนวัดความหนาแน่น} (\texttt{occupancyPolygon}) เป็นรูปหลายเหลี่ยมอีกรูปหนึ่งที่ครอบพื้นที่ถนนในช่วงที่ยาวกว่าโซนนับ ใช้สำหรับนับจำนวนยานพาหนะที่ปรากฏอยู่ในพื้นที่ ณ แต่ละเฟรม โดยไม่สนใจว่ายานพาหนะนั้นข้ามเส้นนับแล้วหรือยัง  เหตุผลที่แยกรูปหลายเหลี่ยมนี้ออกจากโซนนับคือ ทั้งสองค่าตอบคำถามคนละข้อ โซนนับตอบว่ามีรถผ่านไปกี่คัน ส่วนโซนวัดความหนาแน่นตอบว่าขณะนี้มีรถอยู่บนถนนมากน้อยเพียงใด ซึ่งเป็นสองปริมาณที่ต้องใช้ร่วมกันตามที่จะอธิบายในหัวข้อที่ \ref{sec:traffic-classification}
    \item \textbf{ทิศทางที่คาดหวัง} (\texttt{expectedDirection}) เป็นค่าที่ระบุว่าโซนนั้นควรมียานพาหนะเคลื่อนที่ไปในทิศทางใด ยานพาหนะที่ข้ามเส้นในทิศทางตรงข้ามกับค่าที่กำหนดจะถูกปฏิเสธ พร้อมบันทึกเหตุผลไว้ตามที่อธิบายในหัวข้อที่ \ref{subsec:rejection-log}
\end{itemize}

พิกัดของรูปหลายเหลี่ยมและเส้นนับทั้งหมดกำหนดในหน่วยพิกเซลของวิดีโอต้นทาง ผลที่ตามมาคือ ค่าเหล่านี้ผูกกับมุมกล้องและความละเอียดของวิดีโอแต่ละไฟล์ เมื่อเปลี่ยนมุมกล้องหรือเปลี่ยนไฟล์วิดีโอที่มีความละเอียดต่างกัน ต้องกำหนดพิกัดใหม่ทุกครั้ง  ข้อจำกัดนี้ส่งผลต่อไปถึงเกณฑ์ตัดสินสภาพการจราจรด้วย เนื่องจากค่าความหนาแน่นที่วัดได้ขึ้นอยู่กับขนาดของโซนที่วาด จึงต้องปรับเกณฑ์ใหม่ทุกครั้งที่เปลี่ยนมุมกล้อง ตามที่จะกล่าวในหัวข้อที่ \ref{subsec:state-rules}

ในการตั้งค่าที่ใช้ทดสอบ ผู้จัดทำกำหนดโซนตรวจวัดไว้สองโซน ครอบคลุมช่องจราจรทั้งสองฝั่งของถนน โดยโซนหนึ่งรับผิดชอบทิศทางเข้าและอีกโซนหนึ่งรับผิดชอบทิศทางออก เหตุผลที่ต้องแยกสองโซนคือถนนทั้งสองฝั่งไม่ได้ติดขัดพร้อมกัน จากข้อมูลที่วัดได้จริง ครึ่งหนึ่งของหน้าต่างเวลาทั้งหมดมีสถานะของสองฝั่งแตกต่างกัน การรายงานค่าเดียวสำหรับทั้งภาพจึงกลบข้อมูลที่สำคัญไประบบจึงรายงานสถานะแยกรายโซนควบคู่ไปกับค่ารวม โซนนับ เส้นนับ และโซนวัดความหนาแน่น
ที่วาดทับบนเฟรมจริง แสดงในรูปที่ \ref{fig:zone-overlay}

\begin{figure}[ht]
    \centering
    \includegraphics[width=130mm]{fig_3_4.png}
    \caption{โซนนับ เส้นนับ และโซนวัดความหนาแน่นที่กำหนดบนภาพจากกล้อง}
    \label{fig:zone-overlay}
\end{figure}

\subsection{ชุดข้อมูลที่ใช้ทดสอบ}
\label{subsec:datasets}

ผู้จัดทำบันทึกวิดีโอไว้หลายไฟล์เพื่อใช้ในขั้นตอนที่ต่างกัน รายละเอียดแสดงในตารางที่ \ref{tab:datasets}

\begin{table}[ht]
    \caption{รายการวิดีโอที่ใช้ในการพัฒนาและทดสอบระบบ}
    \label{tab:datasets}
    \centering
    \footnotesize
    \begin{tabular}{|l|c|c|c|>{\raggedright\arraybackslash}p{32mm}|>{\raggedright\arraybackslash}p{34mm}|}
        \hline
        \textbf{ชื่อไฟล์} & \textbf{ความละเอียด} & \textbf{\begin{tabular}{@{}c@{}}อัตรา\\เฟรม\end{tabular}} & \textbf{ความยาว} & \multicolumn{1}{c|}{\textbf{สภาพการจราจร}} & \multicolumn{1}{c|}{\textbf{ใช้ในขั้นตอน}} \\
        \hline
        \texttt{traffic-1} & $3840\times2160$ & 30 & 33.7 วินาที & [ระบุ] & พัฒนาและกำหนดโซนเริ่มต้น \\
        \hline
        \texttt{traffic-2} & $3840\times2160$ & 30 & 83.3 วินาที & [ระบุ] & ทดสอบความแม่นยำการนับยานพาหนะ (หัวข้อ 4.3) \\
        \hline
        \texttt{traffic-4} & $3840\times2160$ & 30 & 21.4 วินาที & [ระบุ] & ทดสอบความแม่นยำการนับยานพาหนะ (หัวข้อ 4.3) \\
        \hline
        \texttt{traffic-5} & $3840\times2160$ & 30 & 92.8 วินาที & การจราจรปกติ & วัดช่วงค่าของความหนาแน่นและอัตราการไหล \\
        \hline
        \texttt{traffic-6} & $3840\times2160$ & 30 & 68.3 วินาที & [ระบุ] & ทดสอบความแม่นยำการนับยานพาหนะ (หัวข้อ 4.3) \\
        \hline
        \texttt{event-1} & $1920\times1080$ & 25 & 21.0 วินาที & [ระบุ] & ประมวลผลและทดสอบระบบ \\
        \hline
        \texttt{event-2} & $1920\times1080$ & 60 & 205.9 วินาที & หลากหลาย (ปกติ/เคลื่อนตัวช้า/หยุดนิ่ง) & สาธิตระบบต้นแบบและทดสอบการจำแนกสภาพจราจร (หัวข้อ 4.4) \\
        \hline
        
        \hline
    \end{tabular}
\end{table}

ไฟล์ \texttt{event-1} เป็นคลิปที่ตัดและแปลงมาจากวิดีโอต้นฉบับความละเอียดสูง จึงมีความละเอียดและอัตราเฟรมต่างจากอีกสองไฟล์

ความละเอียดของวิดีโอมีผลโดยตรงต่อความสามารถในการตรวจจับ เนื่องจากขนาดภาพที่ป้อนเข้าแบบจำลองถูกย่อลงก่อนประมวลผล ยานพาหนะขนาดเล็กที่อยู่ไกลจากกล้อง โดยเฉพาะรถจักรยานยนต์ จึงอาจเหลือจำนวนพิกเซลน้อยเกินกว่าที่แบบจำลองจะตรวจจับได้  ประเด็นนี้เป็นสาเหตุหนึ่งของการเลือกค่าขนาดภาพนำเข้าตามที่อธิบายในหัวข้อที่ \ref{subsec:model-params}

ในการวัดช่วงค่าที่เกิดขึ้นจริงจากวิดีโอที่มีการจราจรปกติ ระบบให้ค่าความหนาแน่นอยู่ระหว่าง 0.20 ถึง 2.76 และค่าอัตราการไหลอยู่ระหว่าง 0.0 ถึง 42.4 คันต่อนาที ค่าเหล่านี้ถูกใช้เป็นพื้นฐานในการตั้งเกณฑ์ตัดสินสภาพการจราจรในหัวข้อที่ \ref{subsec:state-rules}  ช่วงค่าดังกล่าวได้จากสภาพการจราจรปกติเท่านั้น ผู้จัดทำจึงใช้คลิป \texttt{event-2} ซึ่งมีช่วงเวลาที่การจราจรเคลื่อนตัวช้าและหยุดนิ่งเกิดขึ้นจริง เป็นข้อมูลสอบเทียบเกณฑ์ในส่วนของสภาวะติดขัด อย่างไรก็ตาม \texttt{event-2} เป็นเพียงคลิปเดียวจากมุมกล้องเดียว และไม่มีช่วงที่การจราจรหนาแน่นแต่ยังไหลได้ดี จึงยังไม่สามารถยืนยันได้ว่าเกณฑ์ที่ตั้งไว้ใช้ได้ทั่วไปกับช่วงเวลาและสภาพการจราจรอื่น ข้อจำกัดนี้ถูกนำไปพิจารณาในการออกแบบการทดลองในหัวข้อที่ \ref{sec:experiment-design}

\section{การตรวจจับและติดตามยานพาหนะ}
\label{sec:vehicle-detection}

ขั้นตอนนี้รับภาพจากวิดีโอทีละเฟรม และคืนรายการยานพาหนะที่ตรวจพบในเฟรมนั้น โดยแต่ละรายการประกอบด้วยรหัสประจำตัวที่คงเดิมข้ามเฟรม ชนิดของยานพาหนะ ค่าความเชื่อมั่นของการตรวจจับ และพิกัดกรอบล้อมในภาพ ผลลัพธ์ชุดนี้ถูกส่งต่อไปยังขั้นตอนการนับในหัวข้อที่ \ref{sec:counting}

การตรวจจับและการติดตามถูกอธิบายรวมกันในหัวข้อเดียว เนื่องจากไลบรารีที่ใช้ดำเนินการทั้งสองอย่างภายในการเรียกใช้ครั้งเดียวกัน การแยกเป็นสองขั้นตอนจึงเป็นการสร้างเส้นแบ่งที่ไม่มีอยู่จริงในการทำงานของระบบ

\subsection{การเลือกแบบจำลองและการกำหนดตัวแปร}
\label{subsec:model-params}

\textbf{แบบจำลองที่ใช้} ผู้จัดทำเลือกใช้ YOLOv8 ขนาด s ซึ่งเป็นขนาดที่สองจากเล็กที่สุดในตระกูลเดียวกัน โดยใช้น้ำหนักที่ฝึกมาจากชุดข้อมูล COCO  เหตุผลของการเลือกขนาดนี้คือความสมดุลระหว่างความแม่นยำกับความเร็วในการประมวลผล ขนาดที่เล็กกว่าให้ความเร็วสูงกว่าแต่ตรวจจับวัตถุขนาดเล็กได้แย่ลง ซึ่งเป็นข้อเสียที่รับไม่ได้สำหรับงานนี้เพราะรถจักรยานยนต์เป็นวัตถุขนาดเล็กและเป็นตัวแปรหลักของระบบ ส่วนขนาดที่ใหญ่กว่าให้ความแม่นยำสูงขึ้นแต่ใช้เวลาประมวลผลต่อเฟรมมากขึ้นจนไม่เหมาะกับการทำงานแบบต่อเนื่อง

\textbf{ขนาดภาพนำเข้า} กำหนดไว้ที่ 1280 พิกเซล ค่านี้ถูกปรับขึ้นจากค่าเดิมที่ 640 พิกเซล หลังจากพบว่าระบบตรวจจับรถจักรยานยนต์ที่อยู่ไกลจากกล้องไม่ได้  สาเหตุคือวิดีโอที่ใช้ในขั้นตอนการพัฒนามีความละเอียด $3840\times2160$ พิกเซล เมื่อย่อภาพลงเหลือ 640 พิกเซลก่อนป้อนเข้าแบบจำลอง ภาพถูกย่อลงหกเท่า รถจักรยานยนต์ซึ่งมีขนาดเล็กและบางกว่ารถยนต์อยู่แล้ว จึงเหลือจำนวนพิกเซลน้อยเกินกว่าที่แบบจำลองจะตรวจจับได้  การเพิ่มขนาดภาพนำเข้าเป็น 1280 พิกเซลลดการย่อลงเหลือสามเท่า ทำให้ระบบตรวจจับรถจักรยานยนต์ได้ โดยแลกกับเวลาประมวลผลต่อเฟรมที่เพิ่มขึ้น

\textbf{ค่าขีดแบ่งความเชื่อมั่น} กำหนดไว้ที่ 0.35 การเลือกค่านี้มาจากข้อสังเกตว่าแบบจำลองให้ค่าความเชื่อมั่นกับรถจักรยานยนต์ต่ำกว่ารถยนต์อย่างเป็นระบบ  จากรายการยานพาหนะ 11 รายการที่ระบบนับได้ในคลิปทดสอบ รถจักรยานยนต์ 8 รายการมีค่าความเชื่อมั่นอยู่ระหว่าง 0.40 ถึง 0.55 ในขณะที่รถยนต์ 3 รายการในคลิปเดียวกันมีค่าอยู่ระหว่าง 0.68 ถึง 0.83  หากตั้งค่าขีดแบ่งไว้สูงกว่านี้ เช่นที่ค่า 0.6 ซึ่งเป็นค่าที่ใช้กันทั่วไป รถจักรยานยนต์เกือบทั้งหมดจะถูกตัดทิ้งตั้งแต่ชั้นการตรวจจับ ซึ่งจะทำให้การแยกอัตราการไหลของรถจักรยานยนต์ในหัวข้อที่ \ref{subsec:motorcycle-split} ไม่มีความหมาย  อย่างไรก็ตาม การตั้งค่าขีดแบ่งต่ำย่อมเพิ่มโอกาสที่ระบบจะตรวจจับสิ่งที่ไม่ใช่ยานพาหนะ ผลกระทบส่วนนี้ถูกจำกัดด้วยตัวกรองเชิงพื้นที่และการโหวตชนิดในหัวข้อที่ \ref{sec:counting}

\textbf{หน่วยประมวลผล} ระบบกำหนดให้ใช้หน่วยประมวลผลกราฟิกเป็นค่าเริ่มต้น และตรวจสอบความพร้อมใช้งานก่อนเริ่มทำงาน หากไม่พบหน่วยประมวลผลกราฟิกที่ใช้ได้ ระบบจะแจ้งเตือนอย่างชัดเจนก่อนจะเปลี่ยนไปใช้หน่วยประมวลผลกลางแทน  การแจ้งเตือนนี้ถูกเพิ่มเข้ามาเพราะการเปลี่ยนไปใช้หน่วยประมวลผลกลางโดยเงียบ ทำให้ผู้พัฒนาเข้าใจผิดว่าโปรแกรมทำงานช้าเพราะตัวโปรแกรมเอง ทั้งที่สาเหตุจริงคือการติดตั้งไลบรารีผิดรุ่น

\subsection{การจำแนกชนิดของยานพาหนะ}
\label{subsec:vehicle-classes}

ระบบกำหนดให้รับเฉพาะยานพาหนะสี่ชนิดจากชุดข้อมูล COCO ได้แก่ รถยนต์ รถจักรยานยนต์ รถโดยสาร และรถบรรทุก โดยส่งรายการหมายเลขชนิดเหล่านี้เข้าไปเป็นตัวกรองตั้งแต่ขั้นการเรียกใช้แบบจำลอง วัตถุชนิดอื่นที่แบบจำลองตรวจพบจะถูกตัดทิ้งตั้งแต่ต้นทาง ไม่ถูกส่งต่อไปยังขั้นตอนการนับ

การกำหนดชนิดเป็นสี่ชนิดนี้มีนัยสำคัญสองประการ

ประการแรก ชนิดของยานพาหนะเป็นตัวแปรที่ระบบใช้ตัดสินใจ ไม่ใช่เพียงข้อมูลประกอบ เนื่องจากอัตราการไหลที่นำไปตัดสินสภาพการจราจรนับเฉพาะรถยนต์ รถโดยสาร และรถบรรทุก การจำแนกชนิดผิดพลาดจึงส่งผลโดยตรงต่อคำตัดสิน ต่างจากระบบตรวจนับทั่วไปที่การจำแนกชนิดผิดพลาดกระทบเพียงรายละเอียดของรายงาน

ประการที่สอง ผู้จัดทำตัดชนิดจักรยานออกจากรายการอย่างจงใจ แม้แบบจำลองจะตรวจจับจักรยานได้ก็ตาม  เหตุผลประการแรกคือจักรยานไม่ใช่ยานยนต์ จึงไม่อยู่ในขอบเขตของการวัดกระแสจราจรตามทฤษฎีในบทที่ \ref{chapter:literature-review}  เหตุผลประการที่สองซึ่งสำคัญกว่าคือ จักรยานมีลักษณะทางภาพใกล้เคียงกับรถจักรยานยนต์มาก โดยเฉพาะเมื่ออยู่ไกลจากกล้องหรืออยู่ในมุมเฉียง หากปล่อยให้ทั้งสองชนิดอยู่ในรายการ แล้วแบบจำลองจำแนกรถจักรยานยนต์เป็นจักรยาน รายการนั้นจะถูกนับรวมเข้าไปในอัตราการไหลที่ใช้ตัดสิน ซึ่งเป็นความผิดพลาดที่ตรงกันข้ามกับสิ่งที่ระบบพยายามหลีกเลี่ยง

อย่างไรก็ตาม การตัดชนิดจักรยานออกก็มีผลข้างเคียงที่ต้องระบุไว้ กล่าวคือ หากแบบจำลองจำแนกรถจักรยานยนต์คันหนึ่งเป็นจักรยาน รายการนั้นจะถูกตัดทิ้งตั้งแต่ต้นทางโดยไม่ทิ้งร่องรอยใด ๆ ไว้ในไฟล์ผลลัพธ์ ทำให้ตรวจสอบย้อนหลังไม่ได้ว่าเกิดขึ้นบ่อยเพียงใด ข้อจำกัดนี้ถูกนำไปพิจารณาในการออกแบบการทดลองวัดความแม่นยำของการนับในหัวข้อที่ \ref{subsec:exp-accuracy}

\subsection{การติดตามยานพาหนะด้วย ByteTrack}
\label{subsec:tracking}

การนับยานพาหนะที่ข้ามเส้นนับจำเป็นต้องรู้ว่ายานพาหนะที่ปรากฏในเฟรมปัจจุบันเป็นคันเดียวกับที่ปรากฏในเฟรมก่อนหน้าหรือไม่ มิฉะนั้นระบบจะนับรถคันเดิมซ้ำทุกเฟรมที่มันอยู่ในภาพ ผู้จัดทำจึงใช้ตัวติดตาม ByteTrack ซึ่งกำหนดรหัสประจำตัวให้ยานพาหนะแต่ละคันและรักษารหัสนั้นไว้ตลอดช่วงที่ยานพาหนะอยู่ในภาพ

เหตุผลที่เลือก ByteTrack แทนตัวติดตามอื่น คือแนวคิดหลักของวิธีนี้ที่ไม่ทิ้งกรอบล้อมที่มีค่าความเชื่อมั่นต่ำ แต่เก็บไว้จับคู่ในรอบที่สอง ตามที่อธิบายไว้ในหัวข้อที่ \ref{subsec:bytetrack}  แนวคิดนี้สอดคล้องกับลักษณะข้อมูลของโครงงานนี้โดยตรง เนื่องจากรถจักรยานยนต์เป็นกลุ่มที่ได้ค่าความเชื่อมั่นต่ำอย่างเป็นระบบตามที่แสดงในหัวข้อที่ \ref{subsec:model-params}  ทางเลือกอีกทางหนึ่งคือตัวติดตาม BoT-SORT ซึ่งสลับรหัสประจำตัวน้อยกว่า แต่ใช้เวลาประมวลผลมากกว่า ผู้จัดทำจึงเลือก ByteTrack เป็นค่าตั้งต้น และเปิดให้สลับไปใช้ BoT-SORT ได้ผ่านไฟล์ตั้งค่า

ในการเรียกใช้ ระบบกำหนดให้ตัวติดตามจดจำสถานะข้ามเฟรม หากไม่กำหนดเช่นนี้ รหัสประจำตัวจะถูกกำหนดใหม่ทุกเฟรม ซึ่งทำให้การนับข้ามเส้นไม่มีความหมาย

ระบบยังกำหนดกฎเพิ่มเติมไว้หนึ่งข้อสำหรับกรณีที่ตัวติดตามยังไม่ได้กำหนดรหัสประจำตัวให้กรอบล้อมในเฟรมนั้น ซึ่งเกิดขึ้นได้จริงในช่วงเฟรมแรก ๆ ของวิดีโอ  ในกรณีนี้ระบบจะข้ามทั้งเฟรมไปโดยไม่ประมวลผลต่อ แทนที่จะกำหนดรหัสประจำตัวขึ้นเอง เหตุผลคือรหัสที่กำหนดขึ้นเองจะไม่ตรงกับรหัสที่ตัวติดตามกำหนดในเฟรมถัดไป ทำให้ระบบเข้าใจว่าเป็นยานพาหนะคนละคันและนับซ้ำ การยอมเสียข้อมูลบางเฟรมจึงปลอดภัยกว่าการเดา

\subsection{ผลลัพธ์ของขั้นตอน}
\label{subsec:detection-output}

เมื่อประมวลผลเฟรมหนึ่งเสร็จ ระบบคืนรายการของยานพาหนะที่ตรวจพบ โดยแต่ละรายการประกอบด้วยรหัสประจำตัวจากตัวติดตาม หมายเลขชนิดของยานพาหนะ ค่าความเชื่อมั่นของการตรวจจับ และพิกัดมุมบนซ้ายกับมุมล่างขวาของกรอบล้อมในหน่วยพิกเซล ข้อมูลชุดนี้เป็นข้อมูลนำเข้าของขั้นตอนการนับในหัวข้อที่ \ref{sec:counting}

\section{การนับยานพาหนะข้ามเส้นและการปฏิเสธการนับ}
\label{sec:counting}

ขั้นตอนนี้รับรายการยานพาหนะที่ตรวจพบในแต่ละเฟรมจากหัวข้อที่ \ref{sec:vehicle-detection} แล้วตัดสินว่ายานพาหนะคันใดข้ามเส้นนับไปแล้วบ้าง ผลลัพธ์มีสองชุด ชุดแรกคือรายการการนับที่ได้รับการยอมรับ ชุดที่สองคือรายการที่ถูกปฏิเสธพร้อมเหตุผล

ตรรกะทั้งหมดในขั้นตอนนี้ถูกแยกออกมาเป็นส่วนที่ไม่ขึ้นกับไลบรารีประมวลผลภาพและไม่อ่านเขียนไฟล์ใด ๆ รับเข้าเป็นตัวเลขและคืนออกเป็นรายการเหตุการณ์เท่านั้น  การออกแบบเช่นนี้ทำให้ข้อผิดพลาดของการนับสามารถเขียนเป็นกรณีทดสอบที่ทำงานเสร็จภายในเวลาไม่กี่วินาที แทนที่จะต้องเปิดวิดีโอดูใหม่ทุกครั้ง ซึ่งเป็นปัญหาสำคัญของโค้ดชุดก่อนหน้า

\subsection{การตรวจการข้ามเส้นด้วยการตัดกันของส่วนของเส้นตรง}
\label{subsec:line-crossing}

\textbf{จุดอ้างอิงของยานพาหนะ} ระบบใช้จุดกึ่งกลางของกรอบล้อมเป็นตัวแทนตำแหน่งของยานพาหนะ มิใช่จุดกึ่งกลางของขอบล่างซึ่งเป็นวิธีที่ใช้กันทั่วไปเพราะใกล้เคียงกับจุดสัมผัสพื้น  เหตุผลคือมุมกล้องของโครงงานนี้มีราวสะพานลอยบังยานพาหนะบางส่วนในบางช่วง ทำให้ขอบล่างของกรอบล้อมกระตุกอย่างรุนแรงระหว่างเฟรม ในขณะที่จุดกึ่งกลางเปลี่ยนแปลงน้อยกว่ามาก

\textbf{วิธีตรวจการข้าม} ระบบไม่ได้ตรวจการข้ามด้วยการเปรียบเทียบพิกัดของยานพาหนะกับตำแหน่งของเส้นโดยตรง แต่พิจารณาการเคลื่อนที่ระหว่างสองเฟรมเป็นส่วนของเส้นตรงเส้นหนึ่ง ซึ่งลากจากจุดอ้างอิงในเฟรมก่อนหน้าไปยังจุดอ้างอิงในเฟรมปัจจุบัน แล้วตรวจว่าส่วนของเส้นตรงนี้ตัดกับส่วนของเส้นตรงที่เป็นเส้นนับหรือไม่ 

การเลือกวิธีนี้แทนการเปรียบเทียบพิกัดมีเหตุผลสามประการ ประการแรก เส้นนับสามารถวางเอียงได้ตามแนวถนน ไม่จำเป็นต้องขนานกับขอบภาพ ประการที่สอง ยานพาหนะที่เคลื่อนที่เร็วอาจเลื่อนตำแหน่งได้เกินสี่สิบพิกเซลต่อเฟรม ซึ่งอาจข้ามเส้นไปทั้งเส้นภายในเฟรมเดียวโดยไม่เคยมีเฟรมใดที่อยู่บนเส้นพอดี ประการที่สาม กรอบล้อมที่สั่นอยู่รอบ ๆ เส้นจะทำให้การเปรียบเทียบพิกัดให้ผลว่าข้ามเส้นซ้ำหลายครั้ง

การตรวจการตัดกันอาศัยเครื่องหมายของผลคูณไขว้ โดยกำหนดให้จุด $a$ และ $b$ เป็นปลายทั้งสองของเส้น และจุด $p$ เป็นจุดที่ต้องการตรวจ ค่า $d$ ตามสมการที่ \ref{eq:side-of-line} จะบอกว่าจุด $p$ อยู่ฝั่งใดของเส้น โดยพิจารณาจากเครื่องหมายของค่า มิใช่ขนาด

\begin{equation}
    d = (b_x - a_x)(p_y - a_y) - (b_y - a_y)(p_x - a_x)
    \label{eq:side-of-line}
\end{equation}

ส่วนของเส้นตรงสองเส้นตัดกันก็ต่อเมื่อปลายทั้งสองของแต่ละเส้นอยู่คนละฝั่งของอีกเส้นหนึ่ง ซึ่งตรวจได้จากการคำนวณค่าตามสมการที่ \ref{eq:side-of-line} รวมสี่ครั้ง

ในการคำนวณ ระบบจัดให้จุดที่ได้ค่า $d$ เท่ากับศูนย์พอดี คือจุดที่อยู่บนเส้น อยู่ในฝั่งเดียวกับค่าลบ ผลที่ได้คือกรณีที่จุดอ้างอิงทับเส้นพอดีจะไม่ถูกนับว่าข้าม  การตัดสินเช่นนี้เป็นการเลือกโดยจงใจ เนื่องจากกรณีดังกล่าวเกิดขึ้นได้ยากมากกับค่าทศนิยม และการนับให้ข้ามมีความเสี่ยงที่จะนับซ้ำมากกว่าประโยชน์ที่ได้รับ

เมื่อพบว่ามีการข้าม ระบบคำนวณพิกัดของจุดตัดเพื่อบันทึกไว้ว่าการข้ามเกิดขึ้น ณ ตำแหน่งใดบนเส้น ข้อมูลนี้ถูกใช้ในขั้นตอนถัดไปเพื่อตรวจว่าการข้ามนั้นอยู่ภายในโซนที่กำหนดหรือไม่

% \begin{figure}[ht]
%     \centering
%     % \includegraphics[width=120mm]{fig_3_5.png}
%     \caption{การตรวจการข้ามเส้นจากการตัดกันของส่วนของเส้นตรง}
%     \label{fig:line-crossing}
% \end{figure}

\subsection{การโหวตและการล็อกชนิดของยานพาหนะ}
\label{subsec:type-voting}

แบบจำลองตรวจจับวัตถุจำแนกชนิดของยานพาหนะใหม่ทุกเฟรม และผลที่ได้ไม่จำเป็นต้องเหมือนกันทุกเฟรม ยานพาหนะคันเดียวกันอาจถูกจำแนกเป็นรถยนต์ในเฟรมหนึ่งและเป็นรถบรรทุกในอีกเฟรมหนึ่ง โดยเฉพาะขณะที่อยู่ไกลจากกล้องหรืออยู่ในมุมที่ไม่ชัดเจน  หากระบบใช้ผลการจำแนกของเฟรมที่เกิดการข้ามเส้นเพียงเฟรมเดียว ชนิดที่บันทึกได้จะไม่เสถียร

ระบบจึงสะสมผลการจำแนกของทุกเฟรมที่พบยานพาหนะคันนั้นไว้ในรูปของการนับคะแนน โดยแต่ละเฟรมลงคะแนนให้ชนิดที่แบบจำลองจำแนกได้หนึ่งคะแนน  อย่างไรก็ตาม เฉพาะเฟรมที่มีค่าความเชื่อมั่นไม่ต่ำกว่าค่าขีดแบ่งเท่านั้นที่มีสิทธิ์ลงคะแนน เฟรมที่ภาพเบลอหรือยานพาหนะอยู่ไกลเกินไปจึงไม่มีส่วนกำหนดชนิด พร้อมกันนี้ระบบยังสะสมผลรวมของค่าความเชื่อมั่นแยกตามชนิดไว้ด้วย

เมื่อยานพาหนะข้ามเส้นและการข้ามนั้นผ่านการตรวจสอบทุกเงื่อนไขแล้ว ระบบจะตัดสินชนิดจากคะแนนเสียงข้างมาก แล้วล็อกชนิดนั้นไว้ ผลการจำแนกในเฟรมถัดจากนั้นจะไม่เปลี่ยนชนิดที่บันทึกไว้อีก  การล็อกมีความจำเป็นเพราะยานพาหนะยังคงปรากฏในภาพต่อไปหลังข้ามเส้นแล้ว หากปล่อยให้ชนิดเปลี่ยนได้ ข้อมูลที่ส่งออกไปแล้วจะไม่ตรงกับข้อมูลที่แสดงบนภาพ

นอกจากนี้ระบบยังคำนวณค่าความเชื่อมั่นเฉลี่ยของเฉพาะเฟรมที่ลงคะแนนให้ชนิดที่ชนะ แล้วบันทึกไว้กับรายการการนับ ค่านี้ใช้สำหรับตรวจสอบย้อนหลังว่ายานพาหนะคันใดได้รับการจำแนกด้วยความมั่นใจต่ำ ซึ่งเป็นกลุ่มที่ควรตรวจสอบก่อนเมื่อผลการนับผิดพลาด

\subsection{การตรวจทิศทางการข้าม}
\label{subsec:direction-check}

การทราบว่ายานพาหนะข้ามเส้นยังไม่เพียงพอ ระบบต้องทราบด้วยว่าข้ามไปในทิศทางใด เนื่องจากถนนมีการจราจรสองทิศทางและแต่ละโซนรับผิดชอบทิศทางเดียว

ระบบพิจารณาทิศทางจากการเปลี่ยนฝั่งของยานพาหนะเทียบกับเส้นนับ โดยเปรียบเทียบค่า $d$ ตามสมการที่ \ref{eq:side-of-line} ของจุดอ้างอิงในเฟรมปัจจุบันกับของจุดอ้างอิงในเฟรมก่อนหน้า หากค่าเพิ่มขึ้น แสดงว่ายานพาหนะเคลื่อนจากฝั่งลบไปยังฝั่งบวก ซึ่งหมายถึงการเคลื่อนที่ลงด้านล่างของภาพ และคือค่า toward ตามนิยามในหัวข้อที่ \ref{subsec:site-survey} คือการเคลื่อนที่เข้าหากล้อง

เพื่อให้ผลการคำนวณไม่ขึ้นกับลำดับที่ผู้ใช้กำหนดปลายทั้งสองของเส้น ระบบจะจัดเรียงปลายเส้นใหม่ให้ชี้จากซ้ายไปขวาเสมอก่อนคำนวณ การจัดเรียงนี้ทำให้ฝั่งที่ให้ค่าบวกเป็นด้านล่างของภาพเสมอ ไม่ว่าผู้ใช้จะคลิกกำหนดจุดใดก่อน  กรณีที่เส้นตั้งเกือบดิ่งซึ่งจะทำให้การจัดเรียงกำกวม ถูกป้องกันไว้ตั้งแต่ขั้นตอนการอ่านค่าตั้งค่า โดยระบบจะปฏิเสธการทำงานพร้อมแจ้งชื่อโซนที่กำหนดเส้นผิด

\subsection{การป้องกันการนับซ้ำและการล้างข้อมูลติดตาม}
\label{subsec:dedup}

ระบบบันทึกรหัสประจำตัวของยานพาหนะทุกคันที่นับไปแล้วไว้ในชุดข้อมูลหนึ่ง ยานพาหนะที่มีรหัสอยู่ในชุดนี้แล้วจะไม่ถูกนับอีก แม้จะข้ามเส้นซ้ำก็ตาม ซึ่งเกิดขึ้นได้จริงเมื่อกรอบล้อมสั่นไปมารอบเส้น

เนื่องจากวิดีโอที่ประมวลผลอาจมีความยาวมาก จำนวนรหัสประจำตัวที่สะสมไว้จะเพิ่มขึ้นเรื่อย ๆ จนใช้หน่วยความจำมากเกินจำเป็น ระบบจึงล้างข้อมูลของยานพาหนะที่ไม่ปรากฏในภาพมานานเป็นระยะ โดยตรวจสอบทุก 300 เฟรม และลบข้อมูลของยานพาหนะที่ไม่ถูกพบมานานเกิน 150 เฟรม

อย่างไรก็ตาม ชุดข้อมูลรหัสที่นับไปแล้วจะไม่ถูกล้างเลย  ข้อยกเว้นนี้มีความจำเป็นเพราะตัวติดตามสามารถนำรหัสประจำตัวเดิมกลับมาใช้กับยานพาหนะคันใหม่ได้ หากระบบล้างชุดข้อมูลนี้ด้วย ยานพาหนะที่ได้รับรหัสซ้ำจะถูกนับใหม่อีกครั้ง

นอกจากนี้ยังมีกรณีที่ระบบไม่สามารถตรวจการข้ามได้ คือเฟรมแรกที่พบยานพาหนะคันหนึ่ง เนื่องจากยังไม่มีจุดอ้างอิงของเฟรมก่อนหน้าให้ลากเป็นส่วนของเส้นตรง ในกรณีนี้ระบบจะบันทึกตำแหน่งไว้และลงคะแนนชนิดเท่านั้น แล้วเริ่มตรวจการข้ามได้ตั้งแต่เฟรมถัดไป

\subsection{การบันทึกรายการที่ถูกปฏิเสธ}
\label{subsec:rejection-log}

การข้ามเส้นที่ตรวจพบไม่ได้ถูกนับทุกครั้ง ระบบตรวจสอบเงื่อนไขสี่ข้อตามลำดับ และหากไม่ผ่านข้อใด จะบันทึกรายการนั้นไว้พร้อมเหตุผลแทนที่จะทิ้งไปเฉย ๆ เงื่อนไขทั้งสี่เรียงตามลำดับการตรวจสอบดังนี้

\begin{enumerate}
    \item \textbf{จุดตัดอยู่นอกโซนนับ} ระบบตรวจว่าจุดที่ยานพาหนะข้ามเส้นอยู่ภายในรูปหลายเหลี่ยมของโซนหรือไม่ โดยใช้วิธียิงรังสีไปทางด้านขวาแล้วนับจำนวนครั้งที่รังสีตัดกับขอบของรูปหลายเหลี่ยม หากได้จำนวนคี่แสดงว่าจุดนั้นอยู่ภายใน  การตรวจข้อนี้ถูกจัดไว้เป็นลำดับแรกเพราะยานพาหนะที่วิ่งอยู่ในโซนอื่นสามารถตัดกับเส้นของโซนนี้ได้โดยไม่เกี่ยวข้องกัน หากตรวจข้ออื่นก่อน รายการเหล่านั้นจะถูกบันทึกด้วยเหตุผลที่ไม่ตรงกับสาเหตุจริง
    \item \textbf{ทิศทางไม่ตรงกับที่โซนกำหนด} ยานพาหนะที่ข้ามเส้นในทิศทางตรงข้ามกับทิศทางที่โซนนั้นรับผิดชอบจะถูกปฏิเสธ
    \item \textbf{ยานพาหนะนี้ถูกนับไปแล้ว} ตามที่อธิบายในหัวข้อที่ \ref{subsec:dedup}
    \item \textbf{ยังไม่มีคะแนนเสียงเพียงพอที่จะระบุชนิด} เกิดขึ้นเมื่อทุกเฟรมที่พบยานพาหนะคันนั้นมีค่าความเชื่อมั่นต่ำกว่าค่าขีดแบ่ง จึงไม่มีเฟรมใดมีสิทธิ์ลงคะแนน ระบบไม่นับยานพาหนะที่ระบุชนิดไม่ได้ เพราะชนิดเป็นตัวแปรที่ใช้ตัดสินสภาพการจราจรตามที่กล่าวในหัวข้อที่ \ref{subsec:vehicle-classes}
\end{enumerate}

รายการที่ถูกปฏิเสธทั้งหมดถูกบันทึกลงไฟล์แยกพร้อมเหตุผล ลำดับเฟรม รหัสประจำตัว ชนิด โซน ทิศทาง และพิกัดของจุดตัด  บันทึกชุดนี้ทำหน้าที่เป็นร่องรอยการตรวจสอบ ซึ่งมีความสำคัญมากกว่าที่เห็นในครั้งแรก เพราะเมื่อยอดนับต่ำกว่าที่ควรจะเป็น ผู้พัฒนาสามารถตรวจได้ทันทีว่ายานพาหนะหายไปที่ขั้นตอนใด  หากพบรายการถูกปฏิเสธจำนวนมากด้วยเหตุผลใดเหตุผลหนึ่ง แสดงว่าปัญหาอยู่ที่ตรรกะการนับหรือการวางโซน แต่หากบันทึกชุดนี้ว่างเปล่าทั้งที่ยอดนับต่ำ แสดงว่ายานพาหนะหายไปตั้งแต่ชั้นการตรวจจับ ซึ่งเป็นคนละปัญหาและต้องแก้คนละที่  กรณีหลังนี้เกิดขึ้นจริงระหว่างการพัฒนา และเป็นสิ่งที่นำไปสู่การปรับขนาดภาพนำเข้าตามที่กล่าวในหัวข้อที่ \ref{subsec:model-params}

\section{การจำแนกสภาพจราจร}
\label{sec:traffic-classification}

ขั้นตอนนี้เป็นหัวใจของระบบ ทำหน้าที่แปลงผลการตรวจจับรายเฟรมและรายการการนับจากหัวข้อที่ \ref{sec:counting} ให้กลายเป็นคำตัดสินว่าถนนช่วงนั้นกำลังอยู่ในสภาพใด

เหตุผลที่ต้องมีขั้นตอนนี้แยกออกมา สืบเนื่องมาจากข้อจำกัดที่กล่าวไว้ในหัวข้อที่ \ref{subsec:traffic-relationship} กล่าวคือ การนับยานพาหนะที่ข้ามเส้นให้ค่าอัตราการไหลเพียงอย่างเดียว และอัตราการไหลเพียงลำพังไม่สามารถแยกสภาวะที่ถนนว่างออกจากสภาวะที่รถติดจนหยุดนิ่งได้ เพราะทั้งสองกรณีให้ค่าเป็นศูนย์เท่ากัน  ระบบจึงต้องวัดปริมาณที่สองเพิ่มเข้ามาเพื่อแยกสองสภาวะนี้ออกจากกัน

โดยหลักการ ปริมาณที่ควรวัดเพิ่มคือความเร็วเฉลี่ยหรือความหนาแน่น ผู้จัดทำเลือกวัดความหนาแน่น เนื่องจากการวัดความเร็วจากภาพจำเป็นต้องแปลงพิกัดในภาพให้เป็นระยะทางจริงบนพื้นถนนเสียก่อน ซึ่งผู้จัดทำได้ทดลองทำแล้วพบว่าค่าความเร็วที่คำนวณได้เกินกว่าค่าที่ควรจะเป็นจริงไปมาก ซึ่งเป็นค่าที่เป็นไปไม่ได้บนถนนที่ศึกษา จึงไม่สามารถนำไปใช้ได้ ในทางกลับกัน ข้อมูลที่จำเป็นสำหรับการวัดความหนาแน่นมีอยู่แล้วในผลการตรวจจับทุกเฟรม เพียงแต่ขั้นตอนการนับไม่ได้นำไปใช้

หลักการสำคัญอีกประการหนึ่งของขั้นตอนนี้คือ \textbf{การแยกการวัดออกจากการตัดสิน} การวัดค่ามีต้นทุนสูงเพราะต้องประมวลผลแบบจำลองตรวจจับวัตถุตลอดทั้งวิดีโอ ในขณะที่การตัดสินมีต้นทุนต่ำมากแต่ต้องปรับเกณฑ์บ่อยครั้ง ระบบจึงบันทึกผลการวัดลงไฟล์ก่อนโดยยังไม่ผนวกคำตัดสินเข้าไป แล้วจึงนำไฟล์นั้นมาตัดสินในภายหลัง  ผลที่ได้คือเมื่อต้องการปรับเกณฑ์ใหม่ ผู้จัดทำไม่จำเป็นต้องประมวลผลวิดีโอซ้ำอีก และตัวเลขที่วัดได้ยังคงเป็นข้อเท็จจริงของวิดีโอที่ไม่เปลี่ยนแปลงตามเกณฑ์ที่เลือกใช้

\subsection{การวัดความหนาแน่น}
\label{subsec:occupancy}

ระบบวัดความหนาแน่นโดยการนับจำนวนยานพาหนะที่ปรากฏอยู่ภายในโซนวัดความหนาแน่นตามที่กำหนดไว้ในหัวข้อที่ \ref{subsec:zone-setup} ในทุกเฟรม แล้วหาค่าเฉลี่ยตลอดช่วงเวลาที่พิจารณา ดังสมการที่ \ref{eq:occupancy}

\begin{equation}
    O = \frac{1}{N}\sum_{i=1}^{N} n_i
    \label{eq:occupancy}
\end{equation}

เมื่อ $O$ คือค่าความหนาแน่นของช่วงเวลานั้น $N$ คือจำนวนเฟรมทั้งหมดในช่วงเวลา และ $n_i$ คือจำนวนยานพาหนะที่อยู่ภายในโซนวัดความหนาแน่นในเฟรมที่ $i$ หน่วยของค่านี้คือคันต่อเฟรม

ความแตกต่างที่สำคัญที่สุดระหว่างการวัดนี้กับการนับข้ามเส้นคือ การวัดความหนาแน่นไม่สนใจว่ายานพาหนะคันนั้นข้ามเส้นนับไปแล้วหรือยัง ยานพาหนะที่จอดนิ่งอยู่กับที่ก็ถูกนับด้วย  คุณสมบัติข้อนี้เองคือสิ่งเดียวที่ทำให้ระบบแยกสภาวะรถติดหยุดนิ่งออกจากสภาวะถนนว่างได้ เพราะในกรณีแรกจะมียานพาหนะจำนวนมากอยู่ในโซนแม้ไม่มีคันใดข้ามเส้นเลย ส่วนกรณีหลังจะไม่มียานพาหนะอยู่ในโซนเลย

การนับจำนวนยานพาหนะในโซนใช้วิธีตรวจว่าจุดอ้างอิงของยานพาหนะแต่ละคันอยู่ภายในรูปหลายเหลี่ยมหรือไม่ ด้วยวิธียิงรังสีเช่นเดียวกับที่ใช้ในหัวข้อที่ \ref{subsec:rejection-log} และนับเฉพาะยานพาหนะสี่ชนิดที่ระบบสนใจตามที่กำหนดไว้ในหัวข้อที่ \ref{subsec:vehicle-classes}  หากโซนใดไม่ได้กำหนดโซนวัดความหนาแน่นแยกไว้ ระบบจะใช้โซนนับแทนโดยอัตโนมัติ

ข้อจำกัดสำคัญที่ต้องระบุไว้คือ ค่าที่ได้นี้เป็น\textbf{ค่าสัมพัทธ์} มิใช่ความหนาแน่นในหน่วยคันต่อกิโลเมตรตามนิยามในหัวข้อที่ \ref{subsec:traffic-parameters} เนื่องจากระบบไม่ทราบระยะทางจริงบนพื้นถนนที่โซนครอบคลุมอยู่  ค่าที่วัดได้จึงขึ้นอยู่กับขนาดของรูปหลายเหลี่ยมที่ผู้ใช้วาด โซนที่วาดกว้างกว่าย่อมให้ค่าสูงกว่าโซนที่วาดแคบกว่าแม้สภาพการจราจรจะเหมือนกัน ผลที่ตามมาคือเกณฑ์ตัดสินที่ตั้งจากโซนหนึ่งไม่สามารถนำไปใช้กับโซนที่วาดต่างออกไปได้โดยตรง

\subsection{การวัดอัตราการไหลต่อหน้าต่างเวลา}
\label{subsec:flow-rate}

ระบบแบ่งวิดีโอออกเป็นหน้าต่างเวลาที่มีความยาวเท่ากัน ค่าที่ใช้ในการทดสอบคือ 20 วินาทีต่อหน้าต่าง
แล้วนับจำนวนยานพาหนะที่ข้ามเส้นนับภายในแต่ละหน้าต่าง จากนั้นแปลงเป็นอัตราต่อนาที
ดังสมการที่\ref{eq:flow-rate}

\begin{equation}
    V = \frac{c \times 60}{T}
    \label{eq:flow-rate}
\end{equation}

เมื่อ $V$ คืออัตราการไหลในหน่วยคันต่อนาที $c$ คือจำนวนยานพาหนะที่ข้ามเส้นนับภายในหน้าต่างนั้น และ $T$ คือความยาวของหน้าต่างในหน่วยวินาที

การเลือกความยาวหน้าต่างเป็นการแลกเปลี่ยนระหว่างความไวกับความเสถียร หน้าต่างที่สั้นเกินไป
ทำให้ค่าอัตราการไหลแกว่งรุนแรงเพราะยานพาหนะเพียงคันเดียวมีน้ำหนักมาก ส่วนหน้าต่างที่ยาว
เกินไปทำให้ระบบตอบสนองต่อการเปลี่ยนแปลงของสภาพจราจรช้า ค่า 20 วินาทีถูกเลือกเพราะที่
ความยาวนี้ ยานพาหนะหนึ่งคันมีน้ำหนักเท่ากับ 3 คันต่อนาที ซึ่งละเอียดเพียงพอที่จะแยกระดับ
ของสภาพจราจรออกจากกันได้ โดยที่ค่าไม่แกว่งจนใช้งานไม่ได้

\subsection{การแยกอัตราการไหลของรถจักรยานยนต์}
\label{subsec:motorcycle-split}

ระบบไม่ได้คำนวณอัตราการไหลเป็นค่าเดียว แต่แยกออกเป็นสองค่าตั้งแต่ขั้นตอนการวัด ค่าแรกนับเฉพาะรถยนต์ รถโดยสาร และรถบรรทุก ส่วนค่าที่สองนับเฉพาะรถจักรยานยนต์  ในการตัดสินสภาพการจราจร ระบบใช้เฉพาะค่าแรกเท่านั้น ส่วนค่าที่สองถูกส่งออกไปเพื่อการรายงานและการตรวจสอบ แต่ไม่มีผลต่อคำตัดสิน

เหตุผลของการแยกสืบเนื่องโดยตรงจากข้อจำกัดที่กล่าวไว้ในหัวข้อที่ \ref{subsec:mixed-traffic} กล่าวคือ รถจักรยานยนต์สามารถแทรกตัวผ่านช่องว่างระหว่างรถที่ติดขัดไปได้ การที่รถจักรยานยนต์ยังคงเคลื่อนที่ผ่านเส้นนับจึงไม่ได้แปลว่าถนนยังไหลอยู่  หากนับรวมยานพาหนะทุกประเภทเป็นค่าเดียว อัตราการไหลที่วัดได้จะมีค่าสูงอยู่เสมอแม้ในขณะที่ยานพาหนะประเภทอื่นหยุดนิ่งโดยสมบูรณ์ ซึ่งจะกลบสัญญาณของสภาวะรถติดจนระบบตรวจไม่พบ

ผลการประมวลผลจริงยืนยันข้อโต้แย้งนี้อย่างชัดเจน ในหน้าต่างเวลาช่วงวินาทีที่ 10 ถึง 20 ของวิดีโอทดสอบ ระบบวัดค่าความหนาแน่นได้ 16.85 คันต่อเฟรม อัตราการไหลของรถยนต์ รถโดยสาร และรถบรรทุกรวมกันได้ \textbf{0.0} คันต่อนาที ในขณะที่อัตราการไหลของรถจักรยานยนต์ได้ 24.0 คันต่อนาที  กล่าวคือตลอดช่วงเวลา 10 วินาทีนั้น ไม่มีรถยนต์ผ่านเส้นนับแม้แต่คันเดียว แต่มีรถจักรยานยนต์ผ่านไป 4 คัน  หากระบบนับรวมกัน อัตราการไหลที่ได้จะเป็น 24.0 คันต่อนาที ซึ่งอยู่ในระดับที่ระบบจะตัดสินว่าถนนยังเคลื่อนตัวได้ตามปกติ ทั้งที่ความจริงคือรถยนต์ไม่ขยับเลยแม้แต่คันเดียว

\subsection{ตรรกะการตัดสินสี่สถานะ}
\label{subsec:state-rules}

เมื่อได้ค่าความหนาแน่นและอัตราการไหลของหน้าต่างเวลาหนึ่งแล้ว ระบบตัดสินสถานะของโซนนั้นจากการพิจารณาสองค่าประกอบกัน โดยตรวจเงื่อนไขตามลำดับที่กำหนด ดังแสดงในตารางที่ \ref{tab:state-rules}

\begin{table}[ht]
    \caption{เงื่อนไขการตัดสินสภาพการจราจรทั้งสี่สถานะ}
    \label{tab:state-rules}
    \centering
    \footnotesize
    \begin{tabular}{|c|l|>{\raggedright\arraybackslash}p{40mm}|>{\raggedright\arraybackslash}p{48mm}|}
        \hline
        \textbf{ลำดับ} & \textbf{สถานะ} & \multicolumn{1}{c|}{\textbf{เงื่อนไข}} & \multicolumn{1}{c|}{\textbf{ความหมาย}} \\
        \hline
        1 & \texttt{normal} & $O < 4.0$ & ยานพาหนะในโซนมีจำนวนน้อย การจราจรเป็นปกติ \\
        \hline
        2 & \texttt{standstill} & $O \geq 4.0$ และ $V_c < 7.5$ & ยานพาหนะหนาแน่นแต่แทบไม่มีคันใดผ่านเส้นนับ คือติดขัดจนหยุดนิ่ง \\
        \hline
        3 & \texttt{slow\_moving} & $O \geq 4.0$ และ $7.5 \leq V_c < 32.0$ & ยานพาหนะหนาแน่นและเคลื่อนตัวได้ช้า \\
        \hline
        4 & \texttt{high\_density} & $O \geq 4.0$ และ $V_c \geq 32.0$ & ยานพาหนะหนาแน่นแต่ยังไหลได้ดี \\
        \hline
    \end{tabular}

    \vspace{4pt}
    \begin{flushleft}
    {\footnotesize หมายเหตุ: $O$ คือค่าความหนาแน่นในหน่วยคันต่อเฟรม และ $V_c$ คืออัตราการไหลของรถยนต์ รถโดยสาร และรถบรรทุก ในหน่วยคันต่อนาที โดยไม่รวมรถจักรยานยนต์ ค่าเกณฑ์ทั้งสามสอบเทียบกับวิดีโอ \texttt{event-2} ที่ใช้หน้าต่างเวลา 20 วินาที}
    \end{flushleft}
\end{table}

\textbf{ลำดับการตรวจเงื่อนไขมีความสำคัญ} ระบบต้องตรวจค่าความหนาแน่นก่อนเสมอ หากตรวจอัตราการไหลก่อน สภาวะถนนว่างซึ่งมีอัตราการไหลเป็นศูนย์เช่นกันจะถูกตัดสินเป็นสภาวะรถติดหยุดนิ่งทันที ซึ่งเป็นความผิดพลาดชนิดเดียวกับที่หัวข้อที่ \ref{subsec:traffic-relationship} ได้ชี้ให้เห็นไว้

ในกรณีที่ระบบกำหนดโซนไว้มากกว่าหนึ่งโซน ระบบจะตัดสินสถานะของแต่ละโซนแยกกันก่อน แล้วเลือกสถานะที่รุนแรงที่สุดมาเป็นตัวแทนของทั้งภาพ โดยเรียงลำดับความรุนแรงจากน้อยไปมากคือ การจราจรปกติ การจราจรหนาแน่น การเคลื่อนตัวช้า และการหยุดนิ่ง  เหตุผลคือการที่ถนนฝั่งใดฝั่งหนึ่งติดขัดก็ถือว่าพื้นที่นั้นมีปัญหาการจราจรแล้ว ไม่ควรถูกกลบด้วยฝั่งที่ยังไหลได้ดี

ค่าเกณฑ์ทั้งสามที่ปรากฏในตารางถูกกำหนดจากช่วงค่าที่วัดได้จริงตามที่รายงานไว้ในหัวข้อที่ \ref{subsec:datasets}  อย่างไรก็ตาม ต้องระบุข้อจำกัดไว้อย่างชัดเจนสามประการ  ประการแรก เกณฑ์ในส่วนของสภาวะติดขัดสอบเทียบกับวิดีโอ \texttt{event-2} เพียงคลิปเดียว ซึ่งมีช่วงเวลาที่รถติดขัดไม่มากนักและไม่มีสถานะการจราจรหนาแน่นปรากฏ เกณฑ์ $V_c \geq 32.0$ ของสถานะ \texttt{high\_density} จึงยังไม่ได้รับการยืนยันด้วยข้อมูลจริงประการที่สอง ค่าเกณฑ์ความหนาแน่นผูกกับขนาดของโซนที่วาด จึงต้องปรับใหม่ทุกครั้งที่เปลี่ยนมุมกล้องหรือขนาดโซน  ประการที่สาม สถานะทั้งสี่นี้\textbf{มิใช่ระดับการให้บริการ}ตามที่นิยามไว้ในหัวข้อที่ \ref{subsec:los} เนื่องจากการประเมินระดับการให้บริการต้องอาศัยความหนาแน่นในหน่วยหน่วยเทียบเท่ารถยนต์นั่งส่วนบุคคลต่อกิโลเมตรต่อช่องจราจร ซึ่งต้องทราบค่าเทียบเท่า จำนวนช่องจราจร และระยะทางจริงบนพื้นถนน ทั้งสามอย่างนี้ระบบไม่มี  การเปรียบเทียบผลของระบบกับระดับการให้บริการโดยตรงจึงทำไม่ได้

\subsection{การจัดการหน้าต่างเวลาที่ไม่สมบูรณ์}
\label{subsec:partial-window}

เมื่อวิดีโอสิ้นสุดลง หน้าต่างเวลาสุดท้ายมักมีความยาวไม่ครบตามที่กำหนด ซึ่งก่อให้เกิดปัญหากับการคำนวณอัตราการไหล เนื่องจากสมการที่ \ref{eq:flow-rate} ใช้ความยาวของหน้าต่างเป็นตัวหาร ยิ่งหน้าต่างสั้นเท่าใด ยานพาหนะเพียงคันเดียวก็ยิ่งทำให้อัตราการไหลที่คำนวณได้มีค่าสูงผิดปกติมากเท่านั้น  ในทางปฏิบัติ หน้าต่างที่ยาวเพียง 1 วินาทีและมียานพาหนะผ่านไปหนึ่งคัน จะให้อัตราการไหลถึง 60 คันต่อนาที ซึ่งไม่สะท้อนสภาพจริงแต่อย่างใด

ระบบจึงกำหนดกฎไว้ว่า หน้าต่างสุดท้ายที่มีความยาวน้อยกว่าครึ่งหนึ่งของความยาวที่กำหนดไว้จะถูกทิ้งไป ไม่นำไปคำนวณและไม่ส่งออก  การยอมเสียข้อมูลช่วงท้ายไม่กี่วินาทีเป็นการแลกที่คุ้มค่ากว่าการส่งตัวเลขที่ทราบอยู่แล้วว่าไม่น่าเชื่อถือออกไปให้โมดูลปลายทางนำไปแสดงผล

นอกจากนี้ ระบบยังตรวจสอบว่าหน้าต่างนั้นมีเฟรมที่ประมวลผลแล้วอย่างน้อยหนึ่งเฟรมหรือไม่ก่อนจะปิดหน้าต่าง หากไม่มีเฟรมใดเลย หน้าต่างนั้นจะถูกทิ้งเช่นกัน เนื่องจากการหารด้วยจำนวนเฟรมที่เป็นศูนย์ในสมการที่ \ref{eq:occupancy} ไม่มีความหมาย

\section{การส่งข้อมูลและการเชื่อมโยงแบบต่อเนื่อง}
\label{sec:realtime-delivery}

โมดูลระบบหลังบ้านทำหน้าที่เป็นตัวกลางระหว่างโมดูลปัญญาประดิษฐ์กับโมดูลแสดงผล โดยมีหน้าที่หลักสามอย่างคือ ตรวจสอบความถูกต้องของข้อมูลที่รับเข้ามา กระจายข้อมูลต่อไปยังผู้รับทุกราย และเก็บยอดนับสะสมเพื่อให้เรียกดูได้ระหว่างการทำงาน

โมดูลนี้พัฒนาด้วย Node.js ร่วมกับเฟรมเวิร์ก Express รุ่น 5 สำหรับให้บริการผ่าน HTTP และไลบรารี Socket.IO รุ่น 4 สำหรับการสื่อสารผ่านเว็บซ็อกเก็ต เขียนด้วยภาษา TypeScript ในโหมดตรวจสอบชนิดข้อมูลอย่างเข้มงวด ร่วมกับการบังคับตรวจสอบค่าที่เข้าถึงผ่านดัชนีว่าอาจไม่มีอยู่จริง  เหตุผลของการเลือกภาษานี้คือ ข้อมูลที่ไหลผ่านโมดูลมีโครงสร้างที่กำหนดไว้ชัดเจน การประกาศชนิดข้อมูลจึงทำให้ความไม่ตรงกันระหว่างสิ่งที่ส่งมากับสิ่งที่โมดูลคาดหวังถูกตรวจพบตั้งแต่ตอนแปลโปรแกรม แทนที่จะพบตอนโปรแกรมทำงานอยู่

จุดปลายทางที่โมดูลเปิดให้บริการแสดงในตารางที่ \ref{tab:endpoints}

\begin{table}[ht]
    \caption{จุดปลายทางของโมดูลระบบหลังบ้าน}
    \label{tab:endpoints}
    \centering
    \footnotesize
    \begin{tabular}{|l|>{\raggedright\arraybackslash}p{92mm}|}
        \hline
        \textbf{ชนิดคำขอและเส้นทาง} & \multicolumn{1}{c|}{\textbf{หน้าที่}} \\
        \hline
        \texttt{POST /api/ingest} & รับเหตุการณ์การข้ามเส้นนับหนึ่งรายการ ตรวจสอบ ปรับปรุงตัวนับ แล้วกระจายต่อ \\
        \hline
        \texttt{POST /api/traffic-state} & รับสรุปสภาพการจราจรหนึ่งช่วงเวลา ตรวจสอบ แล้วกระจายต่อโดยไม่แปลงค่า \\
        \hline
        \texttt{GET /api/stats} & คืนค่ายอดนับสะสมในปัจจุบัน \\
        \hline
        \texttt{POST /api/stats/reset} & ล้างยอดนับสะสมทั้งหมด \\
        \hline
        \texttt{POST /api/session/reset} & เริ่มรอบทดสอบใหม่ โดยล้างไฟล์บันทึกข้อมูลที่กระจายออกไปและยอดนับสะสมพร้อมกัน โมดูลปัญญาประดิษฐ์เรียกใช้โดยอัตโนมัติก่อนเริ่มส่งข้อมูลทุกครั้ง \\
        \hline
        \texttt{GET /health} & ตรวจสอบว่าโมดูลยังทำงานอยู่ \\
        \hline
        \texttt{GET /} & หน้าเว็บสำหรับตรวจสอบระหว่างการพัฒนา \\
        \hline
    \end{tabular}
\end{table}

\subsection{การตรวจสอบความถูกต้องของข้อมูล}
\label{subsec:validation}

ข้อมูลทุกรายการที่รับเข้ามาต้องผ่านการตรวจสอบก่อนที่จะถูกนำไปประมวลผลหรือกระจายต่อ หากไม่ผ่าน โมดูลจะตอบกลับด้วยรหัสสถานะ HTTP 400 พร้อมข้อความระบุเหตุผล และบันทึกคำเตือนพร้อมข้อมูลดิบที่ถูกปฏิเสธไว้ในบันทึกการทำงาน โดยตัดความยาวไม่เกิน 500 ตัวอักษรเพื่อไม่ให้บันทึกโตเกินจำเป็น

สำหรับเหตุการณ์การข้ามเส้นนับ เงื่อนไขที่ทำให้ข้อมูลถูกปฏิเสธมีดังนี้ ไม่มีเวลาประทับหรือเวลาประทับไม่อยู่ในรูปแบบมาตรฐาน ไม่มีรหัสประจำกล้องหรือเป็นข้อความว่าง ไม่มีรหัสประจำตัวของยานพาหนะหรือเป็นข้อความว่าง ชนิดของยานพาหนะไม่อยู่ในสี่ชนิดที่กำหนด ทิศทางไม่อยู่ในสองค่าที่กำหนด ค่าความเชื่อมั่นไม่ใช่ตัวเลขหรืออยู่นอกช่วงศูนย์ถึงหนึ่ง และค่าเวลาหรือลำดับเฟรมของวิดีโอที่ส่งมามีค่าติดลบหรือไม่ใช่จำนวนจำกัด

สำหรับสรุปสภาพการจราจร เงื่อนไขที่ทำให้ข้อมูลถูกปฏิเสธมีดังนี้ ไม่มีเวลาประทับหรือเวลาประทับไม่อยู่ในรูปแบบมาตรฐาน ไม่มีรหัสประจำกล้องหรือเป็นข้อความว่าง เวลาเริ่มต้นหรือเวลาสิ้นสุดของช่วงเวลาติดลบหรือไม่ใช่จำนวนจำกัด สภาพการจราจรที่ส่งมาไม่อยู่ในสี่สถานะที่กำหนด ไม่มีข้อมูลโซนหรือมีโซนเป็นศูนย์โซน และค่าความหนาแน่นหรืออัตราการไหลของโซนใดมีค่าติดลบหรือไม่ใช่จำนวนจำกัด

เงื่อนไขหนึ่งที่ควรกล่าวถึงเป็นพิเศษคือ เวลาสิ้นสุดของช่วงเวลาต้องมากกว่าเวลาเริ่มต้นเสมอ ช่วงเวลาที่ปลายมาก่อนต้นแสดงว่าข้อมูลผิดพลาด และหากปล่อยผ่านไป การคำนวณอัตราการไหลซึ่งใช้ความยาวของช่วงเวลาเป็นตัวหารจะให้ผลลัพธ์ติดลบ ซึ่งเป็นค่าที่ไม่มีความหมายทางกายภาพ

เหตุผลที่ต้องมีชั้นตรวจสอบนี้ทั้งที่ข้อมูลมาจากโมดูลที่ผู้จัดทำเขียนเองมีสองประการ ประการแรก ชั้นนี้ทำให้ระบุได้ทันทีว่าความผิดพลาดเกิดขึ้นที่ต้นทาง เพราะข้อมูลที่ผิดจะถูกปฏิเสธพร้อมเหตุผลที่จุดนี้ แทนที่จะไหลผ่านไปจนปรากฏเป็นอาการแปลก ๆ ในโมดูลแสดงผลแล้วต้องไล่ย้อนหาสาเหตุ ประการที่สอง ชั้นนี้ทำให้สามารถแทนที่โมดูลต้นทางด้วยตัวจำลองระหว่างการพัฒนาได้ โดยมั่นใจว่าตัวจำลองส่งข้อมูลในรูปแบบเดียวกับของจริง

\subsection{การกระจายข้อมูลผ่านเว็บซ็อกเก็ต}
\label{subsec:broadcast}

โมดูลระบบหลังบ้านกระจายข้อมูลผ่านเว็บซ็อกเก็ตสามช่องทาง ซึ่งแยกกันตามชนิดของข้อมูลและตามผู้รับ

ช่องทางแรก ส่งเหตุการณ์การข้ามเส้นนับในรูปแบบเดิมทุกประการโดยไม่แก้ไขค่าใด ๆ ผู้รับของช่องทางนี้คือหน้าเว็บสำหรับตรวจสอบระหว่างการพัฒนา ซึ่งต้องเห็นข้อมูลตรงกับที่โมดูลต้นทางส่งมาจริง

ช่องทางที่สอง ส่งเหตุการณ์เดียวกันแต่ห่อและตัดข้อมูลบางส่วนออก สำหรับโมดูลแสดงผลโดยเฉพาะ ข้อมูลที่ถูกตัดออกได้แก่ ตัวระบุรุ่นของโครงสร้างข้อมูล เวลาและลำดับเฟรมในวิดีโอต้นทาง และค่าความเชื่อมั่นของการตรวจจับ  เหตุผลของการตัดคือข้อมูลเหล่านี้มีความหมายเฉพาะภายในกระบวนการประมวลผลวิดีโอ หรือถูกนำไปใช้ตั้งแต่ขั้นตรวจสอบและนับแล้ว จึงไม่มีบทบาทในการสร้างหรือเคลื่อนย้ายยานพาหนะในฉากสามมิติ

ช่องทางที่สาม ส่งสรุปสภาพการจราจรต่อช่วงเวลาโดยไม่แปลงค่า ช่องทางนี้มีความถี่ต่ำกว่าสองช่องทางแรกมาก คือประมาณหนึ่งครั้งต่อยี่สิบวินาที ในขณะที่สองช่องทางแรกส่งทุกครั้งที่มียานพาหนะข้ามเส้นนับ

หลักการสำคัญของโมดูลนี้คือ ต้องไม่แปลงค่าของข้อมูลก่อนส่งต่อ โดยมีข้อยกเว้นเพียงกรณีเดียวคือช่องทางที่สองซึ่งระบุไว้เป็นลายลักษณ์อักษรในเอกสารสัญญาข้อมูลของโครงงาน  หลักการนี้ทำให้ข้อมูลที่โมดูลแสดงผลได้รับตรงกับข้อมูลที่โมดูลปัญญาประดิษฐ์ส่งออกมาเสมอ เมื่อผลลัพธ์ผิดพลาดจึงตรวจสอบย้อนกลับได้ว่าเกิดขึ้นที่โมดูลใด

นอกจากนี้ยังต้องเน้นว่า การตัดสินสภาพการจราจรเกิดขึ้นที่โมดูลปัญญาประดิษฐ์ ไม่ใช่ที่โมดูลระบบหลังบ้าน โมดูลระบบหลังบ้านเพียงตรวจสอบและส่งต่อคำตัดสินที่มาแล้วเท่านั้น  การวางหน้าที่เช่นนี้ทำให้ตรรกะการตัดสินอยู่ที่เดียวกับที่วัดค่า ซึ่งสอดคล้องกับการแยกการวัดออกจากการตัดสินตามที่อธิบายในหัวข้อที่ \ref{sec:traffic-classification}

\subsection{การนับสะสมในหน่วยความจำ}
\label{subsec:stats}

เมื่อเหตุการณ์การข้ามเส้นนับผ่านการตรวจสอบ โมดูลระบบหลังบ้านปรับปรุงตัวนับสะสมที่แยกตามชนิดและทิศทางของยานพาหนะ  ค่าที่เก็บประกอบด้วยยอดรวมแยกตามทิศทางเข้าและออก ยอดรวมทั้งหมด และยอดแยกย่อยของยานพาหนะแต่ละชนิดในแต่ละทิศทาง พร้อมรหัสกล้องต้นทางและเวลาที่เริ่มนับ ค่าเหล่านี้เรียกดูได้ผ่านจุดปลายทางที่กำหนดไว้ และล้างได้เมื่อต้องการเริ่มการทดสอบรอบใหม่

ตัวนับทั้งหมดอยู่ในหน่วยความจำของกระบวนการเท่านั้น และถูกล้างเมื่อเริ่มกระบวนการใหม่ ตามเหตุผลที่อธิบายไว้ในหัวข้อที่ \ref{subsec:data-contract}

ข้อกำหนดสำคัญของส่วนนี้คือ การนับสะสมต้องไม่หน่วงหรือขัดขวางการกระจายข้อมูล เนื่องจากการแสดงผลแบบต่อเนื่องเป็นวัตถุประสงค์หลักของระบบ ในขณะที่ยอดนับสะสมเป็นข้อมูลประกอบที่ผู้ใช้เรียกดูเมื่อต้องการ  หากการนับทำให้การกระจายข้อมูลช้าลง จะเป็นการแลกสิ่งที่สำคัญกว่าไปกับสิ่งที่สำคัญน้อยกว่า

\subsection{การทดสอบโมดูล}
\label{subsec:backend-tests}

เนื่องจากชั้นตรวจสอบความถูกต้องเป็นจุดที่ความผิดพลาดจะส่งผลกระทบมากที่สุด ผู้จัดทำจึงเขียนกรณีทดสอบระดับหน่วยครอบคลุมส่วนนี้โดยเฉพาะ รวมทั้งสิ้น 46 กรณี แบ่งเป็นฟังก์ชันตรวจสอบเหตุการณ์การข้ามเส้นนับ 15 กรณี ฟังก์ชันตรวจสอบสรุปสภาพการจราจร 20 กรณี ฟังก์ชันแปลงข้อมูลสำหรับโมดูลแสดงผล 5 กรณี และส่วนตัวนับสะสม 6 กรณี ทั้งหมดผ่านการทดสอบ

นอกเหนือจากการทดสอบระดับหน่วย ผู้จัดทำยังทดสอบการทำงานร่วมกันของระบบ โดยนำไฟล์ผลลัพธ์ที่โมดูลปัญญาประดิษฐ์เคยประมวลผลไว้จำนวน 28 เหตุการณ์ มาส่งเข้าโมดูลระบบหลังบ้านซ้ำ แล้วเปรียบเทียบยอดนับสะสมที่ได้กับจำนวนในไฟล์ต้นฉบับ ผลปรากฏว่าตรงกันทุกรายการ คือรถยนต์เข้า 5 คันออก 13 คัน รถบรรทุกเข้า 0 คันออก 2 คัน และรถจักรยานยนต์เข้า 7 คันออก 1 คัน  พร้อมกันนี้ได้ยืนยันว่าโปรแกรมที่เชื่อมต่อผ่านเว็บซ็อกเก็ตได้รับเหตุการณ์ครบทุกชนิดยานพาหนะ และยืนยันว่าข้อมูลที่มีชนิดยานพาหนะไม่ถูกต้องถูกปฏิเสธด้วยรหัสสถานะ HTTP 400 จริง


\section{การพัฒนาระบบแสดงผลสามมิติ (Unity Frontend)}
\label{sec:visualization}

โมดูลแสดงผลส่วนหน้า (Frontend) ถูกพัฒนาขึ้นบนเกมเอนจิน Unity เพื่อทำหน้าที่รับข้อมูลจากโมดูลระบบหลังบ้านและจำลองสภาพการจราจรบนแบบจำลองสามมิติ โดยมีกระบวนการทำงานหลัก 6 ขั้นตอน ดังนี้

\begin{enumerate}
    \item \textbf{การเชื่อมต่อและดักรอข้อมูลจากเซิร์ฟเวอร์}
    ระบบเริ่มต้นการทำงานโดยเชื่อมต่อไคลเอนต์ (Client) เข้ากับเซิร์ฟเวอร์ด้วยโปรโตคอลเว็บซ็อกเก็ตผ่านไลบรารี Socket.IO จากนั้นระบบจะเฝ้ารอเหตุการณ์ (Event Listener) จากฝั่งเซิร์ฟเวอร์ ซึ่งเมื่อเกิดเหตุการณ์ขึ้น เซิร์ฟเวอร์จะส่งคืนชุดข้อมูลดิบในรูปแบบสตริง (String) เพื่อนำไปประมวลผลต่อ

    \item \textbf{การแปลงโครงสร้างข้อมูล (Data Parsing)}
    ระบบจะดำเนินการแปลงข้อความ JSON ที่ได้รับให้เป็นออบเจกต์ (Object) ในภาษา C\# โดยข้อมูลจะถูกจับคู่ (Mapping) เข้ากับคลาส \texttt{SpawnVehicleEnvelope} ซึ่งภายในมีโครงสร้างข้อมูลย่อย \texttt{SpawnVehicleData} สำหรับจัดเก็บรายละเอียดของยานพาหนะแต่ละคันอย่างครบถ้วน เช่น รหัสประจำตัว (\texttt{trackId}), ประเภทของยานพาหนะ (\texttt{type}), ทิศทางการเคลื่อนที่ (\texttt{direction}), รหัสกล้องต้นทาง (\texttt{cameraId}) และเวลาที่เกิดเหตุการณ์ (\texttt{timestamp})

    \item \textbf{การจัดคิวข้อมูล (Queue Management)}
    เมื่อกระบวนการแปลงข้อมูลเสร็จสิ้น ระบบจะทำการตรวจสอบความสมบูรณ์ของข้อมูล หากพบว่าข้อมูลไม่เป็นค่าว่าง (Null) จะทำการส่งชุดข้อมูลดังกล่าวเข้าสู่คิวของระบบจัดการหลัก (Manager Script) เพื่อรอการประมวลผลตามลำดับก่อนหลัง (First-In-First-Out)

    \item \textbf{การวิเคราะห์เงื่อนไขและการเลือกเส้นทาง}
    ระบบจัดการจะตรวจสอบคิวการทำงานอย่างต่อเนื่อง เมื่อมีข้อมูลในคิว ระบบจะนำข้อมูลทิศทาง (Inbound หรือ Outbound) มาวิเคราะห์เพื่อกำหนดเส้นทางการเคลื่อนที่ (Route) และวิเคราะห์ประเภทของยานพาหนะ (\texttt{type}) เพื่อจับคู่กับชุดโมเดลสามมิติ (Prefabs) ที่สอดคล้องกัน

    \item \textbf{การสร้างแบบจำลองด้วยพูลออบเจกต์ (Object Pooling)}
    แทนที่จะสร้างวัตถุใหม่ ระบบจะสุ่มเลือกโมเดลรถยนต์จากคลังออบเจกต์ (Object Pool) ที่ถูกสร้างเตรียมไว้ล่วงหน้า จากนั้นจะกำหนดค่าเริ่มต้นต่าง ๆ ได้แก่ พิกัดจุดเกิด หมุนทิศทางให้สอดคล้องกับเส้นทาง กำหนดความเร็วในการเคลื่อนที่ และส่งมอบชุดพิกัดจุดนำทาง (Waypoints) ให้กับแบบจำลองรถยนต์คันดังกล่าว เมื่อเสร็จสิ้น ระบบจะหน่วงเวลา (Delay) ตามระยะเวลาที่กำหนด ก่อนที่จะวนซ้ำกลับไปประมวลผลข้อมูลของยานพาหนะคันถัดไปในคิว

    \item \textbf{การปรับการแสดงผลตามสภาพการจราจร}
    นอกเหนือจากเหตุการณ์ยานพาหนะรายคัน ระบบยังเฝ้ารอเหตุการณ์ \texttt{traffic\_state} ซึ่งส่งคำตัดสินสภาพการจราจรของแต่ละโซนมาทุกหน้าต่างเวลา เมื่อได้รับข้อมูล ระบบจะบันทึกสถานะล่าสุดของโซนนั้นไว้ และใช้กำหนดช่วงความเร็วของยานพาหนะที่สร้างขึ้นในโซนดังกล่าวหลังจากนั้น เช่น โซนที่มีสถานะการจราจรติดขัดหยุดนิ่งหรือเคลื่อนตัวช้า ยานพาหนะที่แสดงจะเคลื่อนที่ช้ากว่าช่วงความเร็วปกติ ทำให้ฉากสามมิติสะท้อนทั้งปริมาณและสภาพการไหลของการจราจรจริง

    \begin{figure}[ht]
        \centering
        \includegraphics[width=140mm]{fig_3_5.png}
        \caption{ตัวอย่างการแสดงผลใน unity}
        \label{fig:Example of the Unity Interface}
    \end{figure}
\end{enumerate}

\section{การออกแบบการทดลองและตัวชี้วัด}
\label{sec:experiment-design}

หัวข้อนี้อธิบายวิธีการที่ใช้ประเมินประสิทธิภาพของระบบที่พัฒนาขึ้น โดยการทดลองถูกแบ่งออกเป็น 4 ชุดหลัก ประกอบด้วยการทดลอง 3 ชุดแรกที่มุ่งเน้นการประเมินความถูกต้องของผลลัพธ์ในแต่ละขั้นตอนการทำงาน และการทดลองชุดสุดท้ายที่ใช้สำหรับประเมินความหน่วงเวลา (Latency) ของการเชื่อมโยงข้อมูล

การออกแบบการทดลองยึดหลักว่า \textbf{เกณฑ์การยอมรับต้องกำหนดไว้ล่วงหน้า} ก่อนที่จะเห็นผลลัพธ์ เพื่อไม่ให้เกณฑ์ถูกปรับตามผลที่ได้ในภายหลัง เกณฑ์ทั้งหมดจึงรวบรวมไว้ในหัวข้อที่ \ref{subsec:acceptance}

\subsection{การทดลองที่ 1 ความแม่นยำของการนับยานพาหนะ}
\label{subsec:exp-accuracy}

\textbf{วัตถุประสงค์} เพื่อตรวจสอบว่าจำนวนยานพาหนะที่ระบบนับได้ตรงกับจำนวนที่เกิดขึ้นจริงในวิดีโอเพียงใด และเพื่อตรวจสอบว่าความแม่นยำของการนับรถจักรยานยนต์ต่างจากยานพาหนะประเภทอื่นมากน้อยเพียงใด

\textbf{วิธีดำเนินการ} ผู้จัดทำเปิดวิดีโอทดสอบและบันทึกยานพาหนะทุกคันที่เคลื่อนที่ผ่านเส้นนับด้วยตาเปล่า โดยบันทึกเวลาในวิดีโอ ชนิด และทิศทางของแต่ละคัน ผลที่ได้ถือเป็นค่าอ้างอิงที่ถูกต้อง จากนั้นนำมาจับคู่กับรายการที่ระบบบันทึกไว้ในไฟล์ผลลัพธ์ โดยถือว่ารายการสองรายการเป็นคันเดียวกันเมื่อเวลาในวิดีโอต่างกันไม่เกินหนึ่งวินาที

\textbf{ตัวชี้วัด} จากการจับคู่ข้างต้น ผู้จัดทำจำแนกผลออกเป็นสามกลุ่ม ได้แก่ รายการที่ระบบนับถูกต้องทั้งชนิดและทิศทาง ($TP$) รายการที่ระบบนับแต่ไม่มีอยู่จริงหรือระบุชนิดหรือทิศทางผิด ($FP$) และรายการที่เกิดขึ้นจริงแต่ระบบไม่ได้นับ ($FN$) แล้วคำนวณตัวชี้วัดตามสมการที่ \ref{eq:precision} ถึง \ref{eq:f1}

\begin{equation}
    \text{Precision} = \frac{TP}{TP + FP}
    \label{eq:precision}
\end{equation}

\begin{equation}
    \text{Recall} = \frac{TP}{TP + FN}
    \label{eq:recall}
\end{equation}

\begin{equation}
    F_1 = \frac{2 \times \text{Precision} \times \text{Recall}}{\text{Precision} + \text{Recall}}
    \label{eq:f1}
\end{equation}

นอกจากนี้ยังรายงานค่าความคลาดเคลื่อนสัมบูรณ์เฉลี่ยของยอดรวมแยกตามชนิดและทิศทาง ตามสมการที่ \ref{eq:mae} เมื่อ $y_j$ คือยอดนับจริงของกลุ่มที่ $j$ $\hat{y}_j$ คือยอดที่ระบบนับได้ และ $k$ คือจำนวนกลุ่มทั้งหมด

\begin{equation}
    MAE = \frac{1}{k}\sum_{j=1}^{k} \left| y_j - \hat{y}_j \right|
    \label{eq:mae}
\end{equation}

\textbf{การรายงานผล} ตัวชี้วัดทั้งหมดรายงานแยกตามชนิดของยานพาหนะ เนื่องจากสมมติฐานของโครงงานนี้คือรถจักรยานยนต์ตรวจจับได้ยากกว่าประเภทอื่น การรายงานเป็นค่ารวมเพียงค่าเดียวจะกลบความแตกต่างที่ต้องการตรวจสอบ

\textbf{ข้อจำกัดที่ทราบล่วงหน้า} ตามที่อธิบายไว้ในหัวข้อที่ \ref{subsec:vehicle-classes} หากแบบจำลองจำแนกรถจักรยานยนต์เป็นจักรยาน รายการนั้นจะถูกตัดทิ้งตั้งแต่ชั้นการตรวจจับโดยไม่ทิ้งร่องรอยไว้ในไฟล์ผลลัพธ์ ความผิดพลาดชนิดนี้จึงปรากฏเป็น $FN$ ในการทดลองนี้ แต่ไม่สามารถแยกออกจาก $FN$ ที่เกิดจากสาเหตุอื่นได้  การจำแนกสาเหตุของ $FN$ แต่ละรายการต้องอาศัยการเปิดดูวิดีโอประกอบเป็นรายกรณี

\subsection{การทดลองที่ 2 ความแม่นยำของการจำแนกสภาพจราจร}
\label{subsec:exp-classification}

\textbf{วัตถุประสงค์} เพื่อตรวจสอบว่าสถานะที่ระบบตัดสินตรงกับสภาพการจราจรที่ปรากฏในวิดีโอเพียงใด

\textbf{วิธีดำเนินการ} ผู้จัดทำแบ่งวิดีโอออกเป็นหน้าต่างเวลาความยาวเท่ากับที่ระบบใช้ แล้วดูวิดีโอแต่ละหน้าต่างและกำกับสถานะที่เห็นด้วยตาเปล่าเป็นหนึ่งในสี่สถานะตามนิยามในตารางที่ \ref{tab:state-rules} ผลที่ได้ถือเป็นค่าอ้างอิง จากนั้นนำมาเปรียบเทียบกับสถานะที่ระบบตัดสิน

\textbf{ตัวชี้วัด} รายงานผลเป็นตารางสับสน (Confusion Matrix) ขนาดสี่คูณสี่ พร้อมค่าความถูกต้องโดยรวม และรายงานค่า Precision และ Recall แยกรายสถานะ

\textbf{ข้อจำกัดที่ทราบล่วงหน้า} การทดลองนี้มีข้อจำกัดสามประการที่ต้องระบุไว้  ประการแรก การกำกับสถานะด้วยตาเปล่าเป็นการตัดสินเชิงอัตวิสัย โดยเฉพาะการแยกระหว่างสถานะการจราจรหนาแน่นกับการเคลื่อนตัวช้า ซึ่งไม่มีเส้นแบ่งที่ชัดเจนในสายตามนุษย์  ประการที่สอง วิดีโอที่มีสภาพการจราจรติดขัดมีเพียงคลิป \texttt{event-2} คลิปเดียว และไม่มีช่วงที่การจราจรหนาแน่นแต่ยังไหลได้ดี ชุดข้อมูลจึงมีสถานะไม่ครบทั้งสี่ประเภท ทำให้ตารางสับสนมีช่องที่ว่างเปล่าและไม่สามารถประเมินสถานะที่ไม่ปรากฏในข้อมูลได้  ประการที่สาม จำนวนหน้าต่างเวลาที่ได้จากวิดีโอที่มีอยู่มีจำนวนน้อย ผลที่ได้จึงเป็นเพียงการตรวจสอบเบื้องต้น มิใช่การประเมินทางสถิติที่มีนัยสำคัญ

การแก้ข้อจำกัดสองข้อหลังจำเป็นต้องบันทึกวิดีโอเพิ่มในช่วงเวลาและสภาพการจราจรที่หลากหลายกว่าเดิม ซึ่งจัดไว้เป็นงานที่ต้องดำเนินการต่อ

\subsection{การทดลองที่ 3 ผลของการแยกอัตราการไหลของรถจักรยานยนต์}
\label{subsec:exp-ablation}

\textbf{วัตถุประสงค์} เพื่อตรวจสอบว่าการแยกอัตราการไหลของรถจักรยานยนต์ออกจากเกณฑ์ตัดสินตามที่อธิบายในหัวข้อที่ \ref{subsec:motorcycle-split} ส่งผลต่อคำตัดสินของระบบจริงหรือไม่ และส่งผลมากน้อยเพียงใด  การทดลองนี้เป็นการตรวจสอบข้อเสนอหลักของโครงงานโดยตรง

\textbf{วิธีดำเนินการ} การทดลองนี้อาศัยคุณสมบัติของระบบที่แยกการวัดออกจากการตัดสิน ผู้จัดทำนำไฟล์ผลการวัดที่บันทึกไว้แล้วมาตัดสินซ้ำสองครั้งด้วยเกณฑ์ชุดเดียวกัน  ครั้งแรกใช้อัตราการไหลที่นับเฉพาะรถยนต์ รถโดยสาร และรถบรรทุก ซึ่งเป็นวิธีที่ระบบใช้จริง  ครั้งที่สองใช้อัตราการไหลที่รวมรถจักรยานยนต์เข้าไปด้วย ซึ่งเป็นวิธีที่ระบบตรวจนับทั่วไปใช้  เนื่องจากทั้งสองครั้งใช้ผลการวัดชุดเดียวกัน ความแตกต่างที่ปรากฏจึงเกิดจากการเปลี่ยนวิธีคิดอัตราการไหลเพียงอย่างเดียว ไม่ปะปนกับความแปรปรวนจากการประมวลผลภาพ

\textbf{ตัวชี้วัด} รายงานจำนวนหน้าต่างเวลาที่คำตัดสินเปลี่ยนไประหว่างสองวิธี พร้อมระบุว่าเปลี่ยนจากสถานะใดเป็นสถานะใด  จุดที่ให้ความสำคัญเป็นพิเศษคือจำนวนหน้าต่างที่วิธีของโครงงานนี้ตัดสินว่าเป็นสภาวะหยุดนิ่ง แต่วิธีที่นับรวมกลับตัดสินว่าเป็นสภาวะอื่น เนื่องจากเป็นกรณีที่แสดงให้เห็นว่าการนับรวมทำให้ระบบมองไม่เห็นสภาวะรถติดขั้นรุนแรง

\subsection{การทดลองที่ 4 ความหน่วงของการเชื่อมโยงข้อมูล}
\label{subsec:exp-latency}

\textbf{วัตถุประสงค์} เพื่อวัดระยะเวลาตั้งแต่ระบบตรวจพบเหตุการณ์จนถึงเวลาที่ผลปรากฏบนแบบจำลองสามมิติ

\textbf{วิธีดำเนินการ} ผู้จัดทำแทรกการบันทึกเวลาไว้ที่จุดสำคัญสี่จุดของเส้นทางข้อมูล ได้แก่ เวลาที่โมดูลปัญญาประดิษฐ์สร้างเหตุการณ์ เวลาที่โมดูลระบบหลังบ้านได้รับคำขอ เวลาที่โมดูลระบบหลังบ้านกระจายข้อมูลออกไป และเวลาที่โมดูลแสดงผลสร้างยานพาหนะขึ้นในฉาก แล้วคำนวณผลต่างระหว่างแต่ละจุดเพื่อแยกให้เห็นว่าความหน่วงเกิดขึ้นที่ช่วงใดมากที่สุด

\textbf{ตัวชี้วัด} รายงานค่ามัธยฐานและค่าเปอร์เซ็นไทล์ที่ 95 ของความหน่วงในแต่ละช่วง และของความหน่วงรวมตั้งแต่ต้นทางถึงปลายทาง

\textbf{ข้อสังเกตเชิงออกแบบสองประการ} ประการแรก ความหน่วงของข้อมูลสภาพการจราจรมีค่าอย่างน้อยเท่ากับความยาวของหน้าต่างเวลา เนื่องจากระบบต้องรอให้หน้าต่างปิดก่อนจึงจะสรุปผลได้ ค่านี้เป็นคุณสมบัติของการออกแบบ มิใช่ข้อบกพร่องของการส่งข้อมูล และต้องรายงานแยกจากความหน่วงของเหตุการณ์รายคัน  ประการที่สอง โมดูลแสดงผลจัดคิวการสร้างยานพาหนะโดยเว้นระยะห่างระหว่างคันไว้ ทำให้สร้างยานพาหนะได้ไม่เกินสองคันต่อวินาที หากโมดูลปัญญาประดิษฐ์ส่งข้อมูลมาเร็วกว่านั้น คิวจะสะสมและความหน่วงจะเพิ่มขึ้นเรื่อย ๆ ตามความยาวของคิว การทดลองนี้จึงต้องรายงานความยาวคิวสูงสุดที่เกิดขึ้นประกอบด้วย

\subsection{เกณฑ์การยอมรับ}
\label{subsec:acceptance}

เกณฑ์การยอมรับของการทดลองทั้งสี่ชุดกำหนดไว้ล่วงหน้าดังตารางที่ \ref{tab:acceptance}

\begin{table}[ht]
    \caption{เกณฑ์การยอมรับผลการทดลอง}
    \label{tab:acceptance}
    \centering
    \footnotesize
    \begin{tabular}{|c|>{\raggedright\arraybackslash}p{46mm}|>{\raggedright\arraybackslash}p{62mm}|}
        \hline
        \textbf{การทดลอง} & \multicolumn{1}{c|}{\textbf{ตัวชี้วัด}} & \multicolumn{1}{c|}{\textbf{เกณฑ์ที่ยอมรับได้}} \\
        \hline
        1 & Recall ของรถยนต์ & ไม่ต่ำกว่า 0.85 \\
        \hline
        1 & Recall ของรถจักรยานยนต์ & ไม่ต่ำกว่า 0.70 \\
        \hline
        1 & Precision ของทุกชนิดรวมกัน & ไม่ต่ำกว่า 0.85 \\
        \hline
        2 & ความถูกต้องโดยรวมของการจำแนกสถานะ & ไม่ต่ำกว่า 0.70 \\
        \hline
        2 & การจำแนกสภาวะหยุดนิ่งเป็นการจราจรปกติ & ต้องไม่เกิดขึ้นเลย \\
        \hline
        3 & จำนวนหน้าต่างที่คำตัดสินเปลี่ยนไป & ต้องมีอย่างน้อยหนึ่งหน้าต่าง \\
        \hline
        4 & ความหน่วงของเหตุการณ์รายคัน ค่ามัธยฐาน & ไม่เกิน 1 วินาที \\
        \hline
    \end{tabular}
\end{table}

เกณฑ์ของรถจักรยานยนต์ถูกกำหนดให้ต่ำกว่าของรถยนต์โดยตั้งใจ เนื่องจากรถจักรยานยนต์เป็นวัตถุขนาดเล็ก ได้รับค่าความเชื่อมั่นต่ำกว่าอย่างเป็นระบบตามที่แสดงไว้ในหัวข้อที่ \ref{subsec:model-params} และมักถูกบดบังในสภาพจราจรหนาแน่น การตั้งเกณฑ์เท่ากับรถยนต์จึงไม่สมเหตุสมผลกับลักษณะของปัญหา

เกณฑ์ในการทดลองที่ 2 ที่ระบุว่าการจำแนกสภาวะหยุดนิ่งเป็นการจราจรปกติต้องไม่เกิดขึ้นเลย เป็นเกณฑ์ที่เข้มกว่าข้ออื่น เนื่องจากเป็นความผิดพลาดที่ส่งผลเสียมากที่สุด การรายงานว่าถนนโล่งทั้งที่กำลังติดขัดย่อมเป็นอันตรายต่อการนำไปใช้ตัดสินใจมากกว่าการรายงานว่าติดขัดทั้งที่โล่ง